\documentclass[a4paper]{article}
% Español (México) con XeLaTeX: fontspec para fuentes Unicode, polyglossia
% para la división silábica y los nombres del documento en la variante mexicana.
\usepackage{fontspec}
\usepackage{polyglossia}
\setdefaultlanguage[variant=mexican]{spanish}
% Los nombres chinos van como «pinyin (original)»: xeCJK da a los caracteres
% Han su propia fuente; sin ella XeLaTeX los omite con un aviso.
% FandolSong no cubre algunos caracteres de nombres (U+73FA); WenQuanYi Zen Hei sí.
\usepackage[AutoFallBack=true]{xeCJK}
\setCJKmainfont{FandolSong-Regular.otf}
\setCJKfallbackfamilyfont{\CJKrmdefault}{WenQuanYi Zen Hei}
% polyglossia no traduce los nombres de listings.
\AtBeginDocument{\providecommand{\lstlistingname}{}\renewcommand{\lstlistingname}{Listado}%
  \providecommand{\lstlistlistingname}{}\renewcommand{\lstlistlistingname}{Índice de listados}}
\usepackage{amsmath,amssymb}
\usepackage{graphicx}
\usepackage[margin=2.35cm]{geometry}
\usepackage[most]{tcolorbox}
\usepackage{etoolbox}
\usepackage{listings}
\usepackage{xcolor}
\usepackage{hyperref}
\usepackage{booktabs,longtable,tabularx,array}
\usepackage{subcaption}
\usepackage{float}
\usepackage{enumitem}
\usepackage{microtype}
\usepackage{fancyhdr}
\usepackage{needspace}

\definecolor{kbBack}{HTML}{F2F6FA}
\definecolor{kbFrame}{HTML}{9DB4CC}
\definecolor{kbTitle}{HTML}{52708F}
\definecolor{ibBack}{HTML}{FBF5E7}
\definecolor{ibFrame}{HTML}{C9A961}
\definecolor{ibTitle}{HTML}{7D5F1E}
\definecolor{wbBack}{HTML}{FCF2F0}
\definecolor{wbFrame}{HTML}{D6A096}
\definecolor{wbTitle}{HTML}{9E4F43}
\definecolor{dbBack}{HTML}{F4F7F3}
\definecolor{dbFrame}{HTML}{AEC2AC}
\definecolor{dbTitle}{HTML}{5F7F64}

\tcbset{softbox/.style={
  enhanced,breakable,boxrule=0.8pt,arc=2pt,left=3mm,right=3mm,
  top=2mm,bottom=2mm,coltitle=white,fonttitle=\bfseries,
  toptitle=0.6mm,bottomtitle=0.6mm
}}
\newtcolorbox{knowledgebox}[1]{softbox,colback=kbBack,colframe=kbFrame,colbacktitle=kbTitle,title={#1}}
\newtcolorbox{importantbox}[1]{softbox,colback=ibBack,colframe=ibFrame,colbacktitle=ibTitle,title={#1}}
\newtcolorbox{warningbox}[1]{softbox,colback=wbBack,colframe=wbFrame,colbacktitle=wbTitle,title={#1}}
\newtcolorbox{dialoguebox}[1]{softbox,colback=dbBack,colframe=dbFrame,colbacktitle=dbTitle,title={#1}}

\lstset{
  language=Python,basicstyle=\ttfamily\small,
  keywordstyle=\color{kbTitle}\bfseries,stringstyle=\color{wbTitle},
  commentstyle=\color{dbTitle},breaklines=true,frame=single,numbers=left,
  numberstyle=\tiny\color{gray},captionpos=b,columns=fullflexible,
  keepspaces=true,showstringspaces=false,extendedchars=true
}
\BeforeBeginEnvironment{lstlisting}{\Needspace{12\baselineskip}}

\hypersetup{colorlinks=true,linkcolor=kbTitle,urlcolor=kbTitle,citecolor=wbTitle}
\setlist{itemsep=0.25em,topsep=0.4em}
\setlength{\parindent}{2em}
\setlength{\parskip}{0.25em}
\newcolumntype{L}[1]{>{\raggedright\arraybackslash}p{#1}}

\providecommand{\notetitle}{Notas del curso: [tema del curso]}
\providecommand{\noteauthors}{Elaboradas a partir de los materiales públicos de Stanford CS25}
\providecommand{\notedate}{Versión reescrita en 2026}
\providecommand{\videochannel}{Stanford Online}
\providecommand{\videopublishdate}{2022}
\providecommand{\videoduration}{[duración del video]}
\providecommand{\videourl}{}
\providecommand{\videocoverpath}{cover.jpg}
\providecommand{\coursesite}{https://web.stanford.edu/class/cs25/}
\providecommand{\repourl}{https://github.com/hqhq1025/ai-course-notes}
\providecommand{\skillversion}{wdkns/wdkns-skills@39f1a04}

\newcommand{\figsource}[1]{\par\vspace{-0.3em}{\footnotesize\color{gray}\emph{Fuente: #1}}\par}
\newcommand{\teachervoice}[1]{\begin{knowledgebox}{Nota de clase}#1\end{knowledgebox}}
\newcommand{\term}[1]{\textbf{#1}}

\pagestyle{fancy}
\fancyhf{}
\fancyfoot[C]{\thepage}
\fancyfoot[R]{\footnotesize\color{gray}\href{\repourl}{AI Course Notes}}
\renewcommand{\headrulewidth}{0pt}
\renewcommand{\footrulewidth}{0pt}

\newcommand{\makecscover}{
\begin{titlepage}
\centering
\vspace*{0.7cm}
{\Large Stanford CS25 · Transformers United\par}
\vspace{0.8cm}
{\huge\bfseries \notetitle\par}
\vspace{0.6cm}
{\large \noteauthors\par}
\vspace{0.25cm}
{\large \notedate\par}
\vspace{0.9cm}
\ifdefempty{\videocoverpath}{}{
  \includegraphics[width=0.86\textwidth,height=0.43\textheight,keepaspectratio]{\videocoverpath}\par
}
\vfill
\begin{tcolorbox}[width=0.92\textwidth,colback=black!2!white,colframe=black!45,boxrule=0.7pt,arc=2pt]
\textbf{Autor o canal del video}: \videochannel\par
\textbf{Fecha de publicación}: \videopublishdate\par
\textbf{Duración del video}: \videoduration\par
\textbf{Sitio del curso}: \href{\coursesite}{\nolinkurl{\coursesite}}\par
\ifdefempty{\videourl}{}{\textbf{Enlace del video}: \href{\videourl}{\nolinkurl{\videourl}}\par}
\textbf{Normas de elaboración}: \skillversion\ + las normas source-first / teacher-voice / visual-QA de este repositorio
\end{tcolorbox}
\end{titlepage}
\tableofcontents
\newpage
}


\renewcommand{\notetitle}{Lecture 17: Biomedical Transformers}
\renewcommand{\noteauthors}{Elaborado a partir de la clase de Vivek Natarajan, los subtítulos oficiales hechos por personas, las diapositivas recuperadas y los artículos asignados}
\renewcommand{\notedate}{CS25 V2 · versión reescrita source-first de 2026}
\renewcommand{\videochannel}{Stanford Online}
\renewcommand{\videopublishdate}{Clase 2023-02-21; publicado 2023-05-25}
\renewcommand{\videoduration}{1:08:09}
\renewcommand{\videourl}{https://www.youtube.com/watch?v=nz7_wg5iOlA}
\renewcommand{\videocoverpath}{cover.jpg}

\newcommand{\lecturefigure}[3]{%
\begin{figure}[H]
  \centering
  \includegraphics[width=0.94\textwidth,height=0.52\textheight,keepaspectratio]{slides-images/#1}
  \caption{#2}
  \figsource{#3}
\end{figure}
}

\let\standardtableofcontents\tableofcontents
\renewcommand{\tableofcontents}{%
  \begingroup
  \small
  \setlength{\parskip}{0pt}
  \standardtableofcontents
  \endgroup
}

\begin{document}
\makecscover

\section{Auditoría de fuentes e hilo conductor de la clase}

Esta clase abarca el lenguaje clínico, las secuencias de proteínas y la regulación génica. A primera vista parece un recorrido por cuatro artículos, pero el hilo conductor es muy unitario: \textbf{primero se escriben los objetos biomédicos como secuencias y después se adaptan las capacidades de representación, atención y generación del Transformer a las restricciones de cada dominio.} Por ello, esta clase no puede limitarse a extraer los resultados de Med-PaLM ni reducir la segunda mitad a «el Transformer también sirve para biología». Hay que preguntar en cada caso cuál es la entrada, cuál es la señal de supervisión, dónde están el contexto largo o el ruido, cómo se verifica la salida del modelo y qué no puede concluirse de los resultados experimentales.

\subsection{Cómo se reconstruye esta clase}

El sitio oficial del curso ofrece la grabación y tres trabajos recomendados, pero no publica un slide deck independiente. Esta clase toma como fuente visual principal la grabación oficial en alta definición completa de Stanford Online: el tamaño de imagen es de 1920 por 1080 píxeles; un muestreo cada dos segundos produjo 2,045 fotogramas, el filtrado por estabilidad y diferencia dejó 103 candidatos de alta exhaustividad y, tras compararlos uno por uno manualmente, se conservaron 84 estados didácticos distintos. Todas las imágenes se volvieron a extraer del video original a ancho completo, sin recortar el contenido del lado derecho; los subtítulos oficiales \texttt{en-US} hechos por personas se usan para recuperar las motivaciones, limitaciones, preguntas y respuestas en clase y transiciones del ponente.

\begin{longtable}{@{}L{0.18\textwidth}L{0.28\textwidth}L{0.42\textwidth}@{}}
\toprule
Capa de evidencia & Material local & Función en los apuntes \\
\midrule
Grabación oficial & Video de 1:08:09 y 84 imágenes didácticas & Determina el orden de la clase, las tablas, las fórmulas, las gráficas de resultados y la evidencia que el ponente mostró realmente \\
Subtítulos oficiales hechos por personas & \texttt{lecture17.en.srt} y timed transcript & Recuperan el «porqué», las correcciones de gráficas, la explicación de riesgos, el impacto de ingeniería y los juicios sobre el futuro \\
Artículos asignados & Med-PaLM, ProtNLM, Enformer & Contrastan la definición de las tareas, los mecanismos de los modelos, los objetos experimentales y los límites de las conclusiones \\
Artículos de contexto & Performer, DeepConsensus, PaLM, Flan-PaLM & Completan las fórmulas y el contexto de sistema que la clase muestra de forma condensada \\
Archivos de auditoría & manifest, coverage, ledger, blueprint & Demuestran que cada nodo visual y cada nodo de teacher voice tiene un destino explícito \\
\bottomrule
\end{longtable}

\begin{warningbox}{Límite histórico: 21 de febrero de 2023}
Esta clase reconstruye el estado del conocimiento del día en que se impartió. Med-PaLM 2, la serie Gemini, los modelos biomédicos multimodales, los eventos regulatorios y los benchmark más recientes, todos posteriores, no se incorporan retroactivamente como evidencia de la clase. Lo que el ponente plantea al final es una visión de investigación de foundation biomedical AI, no un sistema terminado que ya unifique las tareas clínicas, de proteínas y genómicas.
\end{warningbox}

\subsection{Un recorrido en tres etapas: clínica, proteínas, genoma}

La portada solo presenta el tema, pero la página de agenda revela el diseño de la clase: primero responde desde los primeros principios «por qué el Transformer», después entra sucesivamente en clinical, proteins y genomics, y al final discute el futuro. Al leer estas dos páginas conviene notar que el orden de los dominios desciende, en términos generales, por niveles de abstracción: del lenguaje natural que usan médicos y pacientes a las secuencias de aminoácidos, y de ahí a los nucleótidos y las señales regulatorias; pero en cada nivel que se desciende la salida del modelo debe redefinirse, y no pueden heredarse los criterios de evaluación de un chatbot.

\lecturefigure{slide-01-title-transformers-biomedicine.jpg}{Biomedical Transformers: título de la clase y ponente}{Diapositiva V001 recuperada de la grabación oficial; Stanford CS25 Lecture 17}

\lecturefigure{slide-02-agenda.jpg}{Ruta de la clase: aplicaciones clínicas, proteínas, genoma y futuro}{Diapositiva V002 recuperada de la grabación oficial; Stanford CS25 Lecture 17}

\teachervoice{Al iniciar, el ponente dice que su objetivo es «peel back the curtains», para que los estudiantes vean la innovación que está ocurriendo en la intersección entre la IA y la biomedicine. No pretende demostrar que un solo modelo domina todas las tareas, sino que elige casos que atraviesan el biomedical stack, de modo que los mecanismos comunes y las diferencias entre dominios se hagan visibles a la vez.}

\begin{importantbox}{Cinco preguntas que recorren toda la clase}
Para cada caso deben responderse en orden: ¿cómo se construye la secuencia de entrada? ¿Por qué es difícil el contexto? ¿De dónde proviene la señal de supervisión? ¿Cómo verifican la salida las personas o los experimentos? ¿Lo que respaldan los resultados es una mejora en un benchmark, un descubrimiento científico, un apoyo clínico o un despliegue real? Si estas cinco capas se mezclan, las puntuaciones de un modelo se presentan fácilmente, por error, como capacidad médica.
\end{importantbox}

\subsection{Resumen de la sección}

La cadena de evidencia de esta clase se compone de la grabación oficial, los subtítulos hechos por personas y cinco grupos de artículos principales. La estructura de la clase no es una «colección de chatbots médicos», sino una comparación de cómo el Transformer procesa distintas secuencias biomédicas, cómo cambia la complejidad, cómo define la salida y cómo los resultados del modelo se reconectan con la verificación clínica o experimental.
\section{Por qué el Transformer es adecuado para los datos biomédicos}

La sección anterior estableció el mapa del curso; esta sección responde la pregunta común desde la capa de representación. La ventaja más directa del Transformer no es que «entienda medicina», sino que puede establecer dependencias contextuales en secuencias largas y recibir, con un esqueleto de cómputo similar, token discretos, eventos, aminoácidos o nucleótidos. Esta perspectiva unificada es poderosa, pero solo es válida si al mismo tiempo se conserva la semántica de cada dominio.

\subsection{Los cuatro tipos de secuencia no son el mismo lenguaje}

El ponente primero pide a los estudiantes presentes que respondan por qué el Transformer es adecuado para la biomedicine y luego aterriza las respuestas en cuatro objetos concretos. La clinical note es una secuencia de texto escrita por el médico; el electronic medical record (EMR, expediente clínico electrónico) puede verse como una secuencia de eventos formada por las interacciones del paciente con el sistema de salud; la protein es una secuencia de aminoácidos unidos por enlaces peptídicos; y el genome es una secuencia de bases nucleotídicas. Todas tienen un orden, pero poseen «vocabularios», escalas temporales y fuentes de error completamente distintos.

\lecturefigure{slide-03-why-transformers-biomedicine.jpg}{Pregunta de primeros principios: ¿por qué el Transformer es adecuado para la biomedicina?}{Diapositiva V003 recuperada del video oficial; Stanford CS25 Lecture 17}

\lecturefigure{slide-04-clinical-notes-sequences.jpg}{Notas clínicas: secuencias de texto formadas por el lenguaje de los médicos}{Diapositiva V004 recuperada del video oficial; Stanford CS25 Lecture 17}

\begin{knowledgebox}{Lectura de la figura: primero distinguir la «unidad de secuencia»}
En las notas clínicas, la posición suele indicar el orden del texto, pero en un registro de eventos la posición puede representar tiempo, consultas o estudios; en las proteínas, la posición corresponde a un amino acid residue, y en el genoma, a un nucleotide. Lo común es que el contexto cambia el significado de la posición actual; la diferencia es que el modelo debe usar tokenización, codificación posicional, función objetivo y formas de verificación adaptadas al dominio.
\end{knowledgebox}

La diapositiva del expediente clínico electrónico enfatiza «encounters with the medical system»: los diagnósticos, recetas, análisis de laboratorio y consultas del paciente no son solo párrafos de lenguaje natural, sino un flujo de eventos con tiempos, intervalos irregulares y valores faltantes. La diapositiva de proteínas desciende hasta el nivel molecular: la disposición lineal de los aminoácidos se pliega en una estructura tridimensional que influye en los sitios de unión y en la función; por eso, el modelado de secuencias puede aportar representaciones, pero no puede prescindir de la estructura ni de los experimentos.

\lecturefigure{slide-05-electronic-medical-records.jpg}{Expediente clínico electrónico: secuencia de eventos de las interacciones del paciente con el sistema de salud}{Diapositiva V005 recuperada del video oficial; Stanford CS25 Lecture 17}

\lecturefigure{slide-06-protein-sequences.jpg}{Proteínas: secuencias biológicas formadas por aminoácidos unidos}{Diapositiva V006 recuperada del video oficial; Stanford CS25 Lecture 17}

\teachervoice{El ponente llama en broma a las clinical notes «doctor gibberish» y enseguida se corrige a doctor speak. La broma encierra un problema real de diseño: quizá el modelo pueda generar terminología médica, pero si el paciente puede entenderla y si el médico puede verificarla con rapidez son dimensiones que la human evaluation posterior debe medir por separado.}

\subsection{El genoma y «las secuencias están en todas partes»}

La diapositiva del genoma muestra pares de bases nucleotídicas y, a continuación, reúne texto clínico, imágenes, proteínas y DNA en un mismo diagrama general. Aquí, multimodal (multimodal) no significa forzar todos los datos a convertirse en una oración, sino permitir que cada modalidad conserve su propia codificación y luego establecer relaciones mediante una representación común o atención entre modalidades. Por ejemplo, una imagen no es por naturaleza una secuencia unidimensional, pero puede dividirse en patch token; las trayectorias de expresión génica tampoco son lenguaje, pero pueden servir como objetivo de predicción en cada posición.

\lecturefigure{slide-07-genome-sequences.jpg}{Genoma: secuencias largas formadas por pares de bases nucleotídicas}{Diapositiva V007 recuperada del video oficial; Stanford CS25 Lecture 17}

\lecturefigure{slide-08-biomedical-sequences-everywhere.jpg}{En la biomedicina, las secuencias y las modalidades están en todas partes}{Diapositiva V008 recuperada del video oficial; Stanford CS25 Lecture 17}

\begin{longtable}{@{}L{0.16\textwidth}L{0.19\textwidth}L{0.22\textwidth}L{0.31\textwidth}@{}}
\toprule
Término & Unidad de secuencia & Contexto principal & Restricciones del dominio que no se pueden ignorar \\
\midrule
clinical note & subword / word & Antecedentes y problema actual & Abreviaturas, negaciones, tiempo, privacidad, consenso clínico \\
EMR event & Consulta/diagnóstico/fármaco/análisis de laboratorio & Trayectoria longitudinal del paciente & Tiempos irregulares, faltantes, sistemas de codificación y sesgo de selección \\
protein sequence & amino acid & motif locales y relaciones entre residuos distantes & Plegamiento tridimensional, conservación evolutiva, función experimental \\
genome sequence & nucleotide & Elementos reguladores e interacciones de largo alcance & Regiones no codificantes, tipo celular, entorno de la cromatina y ruido de medición \\
multimodal input & text/image/signal token & Relaciones dentro de cada modalidad y entre modalidades & Alineación, sesgo de muestreo, escalas de error distintas \\
\bottomrule
\end{longtable}

Este glosario implica además una conclusión importante de ingeniería de datos: el llamado «preentrenamiento a gran escala» no es lo mismo en los cuatro tipos de datos. La escala del texto clínico está limitada por la privacidad, los hábitos de redacción de cada institución y la desidentificación; los eventos del EMR suelen presentar sesgo de selección por los seguros, el acceso a la atención y los sistemas de codificación; muchas secuencias de las bases de datos de proteínas provienen de especies cercanas, por lo que una partición aleatoria puede hacer que el conjunto de entrenamiento y el de prueba compartan parientes cercanos; y los track genómicos se concentran en unas pocas líneas celulares y protocolos experimentales. Si solo se informa el número de token y no la procedencia, la deduplicación o la partición por familia o por individuo, el modelo puede obtener resultados optimistas al memorizar muestras similares. Por lo tanto, el primer trabajo de ingeniería después de la representación de secuencias debe ser definir la independencia: qué fuga bloquea cada partición, ya sea por paciente, por tiempo, por familia de proteínas o por cromosoma.

\subsection{Esqueleto de cómputo común y límites de aplicabilidad}

La diapositiva general condensa las ventajas en tres términos: sequence modeling, long-range dependencies, scaling. La self-attention estándar construye similitudes por pares para una secuencia de longitud $n$, por lo que puede conectar directamente posiciones lejanas; además, el preentrenamiento permite al modelo aprovechar grandes cantidades de datos sin etiquetar o con etiquetas débiles. Pero las ventajas computacionales y estadísticas no traen automáticamente seguridad clínica, veracidad de la función proteica ni causalidad de las variantes; los cuatro casos que siguen buscan precisamente cubrir estas carencias.

\lecturefigure{slide-09-transformer-biomedical-fit.jpg}{Puntos de coincidencia entre el Transformer y los datos biomédicos}{Diapositiva V009 recuperada del video oficial; Stanford CS25 Lecture 17}

\begin{importantbox}{Representación unificada no equivale a evidencia unificada}
Lo que el Transformer aporta es un lenguaje de cómputo unificado: representación por token, interacción contextual y entrenamiento escalable. Las conclusiones médicas siguen requiriendo un ciclo cerrado de evidencia del dominio: las respuestas clínicas las evalúan médicos y pacientes, las descripciones de proteínas deben contrastarse con bases de datos y experimentos, la corrección de errores de secuenciación debe medirse con read-level accuracy, y las predicciones regulatorias deben compararse con genomic tracks reales y variant assays.
\end{importantbox}

\begin{warningbox}{El error de pensar que «todo es una secuencia»}
Colocar los objetos en una fila no significa que el modelo haya aprendido los invariantes correctos. Los intervalos de tiempo del EMR, la proximidad tridimensional de las proteínas y la dependencia del tipo celular en el DNA pueden quedar ocultos por un orden unidimensional. La serialización es el punto de entrada del modelado, no un sustituto de la estructura del dominio.
\end{warningbox}

\subsection{Resumen de la sección}

El texto clínico, los eventos médicos, las proteínas y el genoma comparten el orden y la dependencia del contexto, pero difieren por completo en la semántica de los token, la longitud, el ruido y la forma de validación. El valor del Transformer está en ofrecer un esqueleto de cómputo transferible; un sistema biomédico de alta calidad, en cambio, debe restituir para cada tipo de datos la estructura del dominio y la evidencia independiente.
\section{Preguntas clínicas: del conocimiento médico a MultiMedQA}

Con la perspectiva de secuencias ya establecida, el curso entra primero en el texto clínico, el terreno más cercano a la forma original de los modelos de lenguaje. La pregunta no es «si el modelo puede recitar hechos médicos», sino si, ante preguntas abiertas, puede organizar el consenso científico, cubrir la información necesaria, evitar omisiones peligrosas y redactar la respuesta de modo que la puedan usar tanto los médicos como los usuarios comunes.

\subsection{Por qué es difícil evaluar el conocimiento clínico}

Después de la diapositiva que separa la sección clínica, el ponente presenta \emph{Large Language Models Encode Clinical Knowledge}. La diapositiva de motivación muestra la magnitud y la complejidad del conocimiento médico: enfermedades, síntomas, estudios, tratamientos y antecedentes individuales dependen unos de otros, y el conocimiento especializado se actualiza de forma continua. Una sola puntuación multiple-choice puede medir la capacidad de elegir en casos acotados, pero no indica si una respuesta larga es exacta, si omite señales de alarma importantes ni si de verdad ayuda al paciente.

\lecturefigure{slide-10-section-clinical.jpg}{Cambio de sección: aplicaciones clínicas}{Diapositiva V010 recuperada de la grabación oficial; Stanford CS25 Lecture 17}

\lecturefigure{slide-11-clinical-knowledge-paper.jpg}{Trabajo asignado: Large Language Models Encode Clinical Knowledge}{Diapositiva V011 recuperada de la grabación oficial; Stanford CS25 Lecture 17}

\lecturefigure{slide-12-clinical-motivation.jpg}{Motivación clínica: el conocimiento médico es amplio, de relaciones complejas y de alto riesgo}{Diapositiva V012 recuperada de la grabación oficial; Stanford CS25 Lecture 17}

\begin{knowledgebox}{Concepto previo: los tres niveles de dificultad de la respuesta a preguntas médicas}
El primer nivel es recordar hechos, por ejemplo el efecto de un medicamento; el segundo es combinar varios hechos dentro del contexto de un caso clínico; el tercero es la comunicación y la gestión del riesgo, que incluye detectar la información faltante, expresar la incertidumbre y recomendar acudir al médico. Las preguntas de opción múltiple cubren sobre todo una parte de los dos primeros niveles; solo las respuestas largas ponen al descubierto el tercero, aunque también son más difíciles de calificar de forma automática.
\end{knowledgebox}

\subsection{De los modelos especializados a los grandes modelos de propósito general}

La diapositiva de la línea de tiempo repasa la evolución de los biomedical language model: los primeros métodos solían entrenar modelos especializados con PubMed o con texto clínico; después, los modelos de lenguaje de propósito general a gran escala comenzaron a abordar los benchmark médicos mediante few-shot o instruction-following. La diapositiva siguiente recuerda, con «What's missing?», que la escala y la fluidez todavía no resuelven la confiabilidad médica, la evaluación de respuestas largas ni las diferencias entre usuarios.

\lecturefigure{slide-13-biomedical-llm-timeline.jpg}{Línea de tiempo de los modelos de lenguaje biomédicos}{Diapositiva V013 recuperada de la grabación oficial; Stanford CS25 Lecture 17}

\lecturefigure{slide-14-clinical-whats-missing.jpg}{¿Qué les falta a los modelos de lenguaje médicos existentes?}{Diapositiva V014 recuperada de la grabación oficial; Stanford CS25 Lecture 17}

Al leer la figura, conviene comparar primero los paradigmas de entrenamiento y no solo el número de parámetros: el domain-specific pretraining modifica las representaciones mediante un corpus especializado; el few-shot prompting modifica principalmente el contexto; el instruction tuning entrena al modelo con un conjunto de tareas para que siga instrucciones; y el instruction prompt tuning de Med-PaLM va más allá y aprende las indicaciones suaves que requieren las tareas médicas. Cada uno interviene, respectivamente, sobre el corpus, los parámetros o las indicaciones, por lo que no pueden agruparse bajo el nombre de «ajuste fino médico».

\begin{warningbox}{Que un modelo grande responda preguntas no significa que sepa ejercer la medicina}
Las preguntas de un benchmark suelen presentar ya un problema cerrado y respuestas candidatas, mientras que la práctica clínica real exige reunir información de forma activa, manejar evidencia contradictoria y asumir las consecuencias. La clase solo analiza evidencia de investigación sobre respuesta a preguntas de conocimiento y apoyo textual; no demuestra que el modelo pueda diagnosticar, recetar ni sustituir al médico por sí mismo.
\end{warningbox}

\subsection{MultiMedQA: reunir referencias de evaluación dispersas en un mismo marco}

La diapositiva Medical question answering divide primero la tarea en exámenes profesionales, preguntas sobre investigación médica y preguntas de salud de los consumidores. A continuación, MultiMedQA reúne varios conjuntos de datos, como MedQA, MedMCQA, PubMedQA y MMLU clinical topics, y añade preguntas reales de consumidores del tipo de HealthSearchQA. El valor de esto no está en producir una puntuación total que lo abarque todo, sino en enfrentar al modelo, a la vez, con opciones cerradas y con respuestas largas abiertas.

\lecturefigure{slide-15-medical-question-answering.jpg}{Respuesta a preguntas médicas: de las preguntas de examen a las preguntas reales de los usuarios}{Diapositiva V015 recuperada de la grabación oficial; Stanford CS25 Lecture 17}

\lecturefigure{slide-16-multimedqa-benchmark.jpg}{MultiMedQA: referencia de evaluación de respuesta a preguntas médicas con varios conjuntos de datos}{Diapositiva V016 recuperada de la grabación oficial; Stanford CS25 Lecture 17}

\lecturefigure{slide-17-multimedqa-examples.jpg}{Ejemplos de preguntas de distintas fuentes en MultiMedQA}{Diapositiva V017 recuperada de la grabación oficial; Stanford CS25 Lecture 17}

\lecturefigure{slide-18-multimedqa-summary.jpg}{Alcance de las tareas y niveles de evaluación de MultiMedQA}{Diapositiva V018 recuperada de la grabación oficial; Stanford CS25 Lecture 17}

\begin{knowledgebox}{Lectura de la figura: no promediar directamente distintos conjuntos de datos como si fueran «capacidad médica»}
MedQA se parece al estilo de los exámenes de licencia médica de Estados Unidos; MedMCQA abarca preguntas de opción múltiple de varias disciplinas médicas; PubMedQA se acerca más a juzgar la evidencia de la literatura; MMLU clinical topics es el subconjunto médico de una evaluación general; y las preguntas de los consumidores no tienen opciones candidatas por naturaleza. Difieren en la fuente del conocimiento que miden, en la forma de la respuesta y en la confiabilidad de la calificación; la vista general debe tomarse como un coverage map, no como la medición precisa de una sola variable latente.
\end{knowledgebox}

Desde el punto de vista del diseño de datos, MultiMedQA crea a la vez una oportunidad y una tensión. La oportunidad es que un mismo modelo base puede recibir muchos tipos de preguntas a través de una interfaz de prompting unificada, lo que permite a los investigadores observar si las capacidades se transfieren entre conjuntos de datos; la tensión es que la label semantics de cada conjunto de datos no es la misma: el «yes/no/maybe» de PubMedQA depende del contexto de la literatura proporcionada, la mejor opción de MedQA depende de las normas del examen y las preguntas de los consumidores a menudo requieren aclarar información en lugar de responder de inmediato. Si durante el entrenamiento todas las respuestas se tratan como una token loss homogénea, el modelo puede aprender la forma y no las condiciones de decisión. Un sistema más prudente conserva la dataset identity, el answer type, el evidence requirement y la abstention policy, e informa los resultados de cada subtarea por separado, en vez de usar un promedio macro que oculte en qué tipo de preguntas de usuarios se concentran los fallos.

\teachervoice{El ponente enfatiza que la medicina es un campo knowledge intensive, pero que «el modelo conoce la respuesta» no basta. Lo que realmente falta es un marco de respuestas largas que puedan evaluar tanto los especialistas como los usuarios comunes, porque a los médicos les importan el consenso científico, las omisiones y los daños, y a los pacientes, además, si pueden entender la respuesta y si esta atiende a su intención.}

\subsection{Resumen de la sección}

La contribución de MultiMedQA es, ante todo, ampliar los tipos de preguntas y el rango de usuarios. Reúne en un mismo marco de investigación las preguntas médicas de opción múltiple, las preguntas sobre investigación y las preguntas largas de los consumidores, y al mismo tiempo pone al descubierto un hecho clave: cuanto más se acerca la tarea a la comunicación real, menos puede depender solo del exact match o de la exactitud en opción múltiple.
\section{Cómo evaluar respuestas largas: la doble perspectiva de expertos y usuarios comunes}

La sección anterior explicó por qué se necesitan preguntas abiertas; esta sección aborda un problema más difícil: ¿cómo se juzga si un texto médico sin una única respuesta estándar es bueno o malo? El curso divide la evaluation en varios ejes y pide a los clinicians y a los laypeople que evalúen, por separado, la parte que cada grupo está más capacitado para juzgar.

\subsection{El marco de evaluación no es una calificación global}

La diapositiva del Evaluation framework es la clave de lectura de todas las figuras de resultados que siguen. La evaluación de expertos se centra en la concordancia con el scientific consensus, en si la respuesta muestra una comprehension correcta o incorrecta, en la retrieval y el reasoning, en si hay missing content, y en el harm y el bias potenciales; la evaluación de usuarios comunes, en cambio, se centra en si la respuesta atiende la intención de la pregunta y si es helpful. Cada eje corresponde a un modo de falla distinto, por lo que no se puede usar un promedio para ocultar una deficiencia peligrosa.

\lecturefigure{slide-19-evaluation-framework.jpg}{Marco de evaluación humana de respuestas largas}{Diapositiva V019 recuperada de la grabación oficial; Stanford CS25 Lecture 17}

\lecturefigure{slide-20-clinician-evaluation-rubric.jpg}{Dimensiones de evaluación de los expertos clínicos}{Diapositiva V020 recuperada de la grabación oficial; Stanford CS25 Lecture 17}

\lecturefigure{slide-21-layperson-evaluation-rubric.jpg}{Dimensiones de evaluación de los usuarios comunes}{Diapositiva V021 recuperada de la grabación oficial; Stanford CS25 Lecture 17}

\begin{longtable}{@{}L{0.22\textwidth}L{0.32\textwidth}L{0.34\textwidth}@{}}
\toprule
Eje de evaluación & Pregunta principal & Falla típica \\
\midrule
scientific consensus & ¿La respuesta concuerda con el conocimiento médico confiable de su momento? & Hace afirmaciones que contradicen el consenso o carecen de evidencia suficiente \\
comprehension / retrieval / reasoning & ¿Entiende la pregunta, recupera el conocimiento pertinente y lo combina correctamente? & Una misma respuesta contiene inferencias correctas e incorrectas \\
incorrect or missing content & ¿Agrega errores irrelevantes u omite contenido esencial? & Detallada pero fuera de tema; pasa por alto señales de alarma que exigen atención médica \\
extent / likelihood of harm & ¿Qué tan grave y qué tan probable es el daño que puede causar el error? & Subestima la enfermedad, induce a retrasar el tratamiento, lleva a un autotratamiento equivocado \\
bias & ¿Aparecen sesgos demográficos o sociales inapropiados? & Vincula sin evidencia atributos de un grupo con un diagnóstico o un juicio de valor \\
intent / helpfulness & ¿Responde a lo que realmente le importa a quien pregunta? & Técnicamente correcta, pero incomprensible o imposible de llevar a la práctica \\
\bottomrule
\end{longtable}

\begin{importantbox}{La perspectiva vectorial de la evaluación médica}
Cada respuesta se representa como un vector de evaluación $\mathbf{e}=(c,r,m,h,b,u)$, cuyas componentes son, respectivamente, consenso, razonamiento, omisión, daño, sesgo y utilidad. Mejorar el sistema no consiste en subir una calificación global, sino en buscar las coordenadas inaceptables: por ejemplo, si la helpfulness mejora pero el harm también aumenta, no puede hablarse de una mejora neta.
\end{importantbox}

Esta evaluación vectorial también cambia la forma de iterar sobre el modelo. Supongamos que un nuevo prompt eleva el scientific consensus, pero al mismo tiempo aumenta el missing content y la verbosity; el equipo tendría que revisar los párrafos añadidos en lugar de declarar que la versión es mejor en conjunto. Si la layperson helpfulness mejora pero la evaluación de harm de los clínicos empeora, la optimización de legibilidad orientada a pacientes rebasó el límite de seguridad. En una auditoría real, cada cambio del modelo puede registrarse en un change log que indique «qué dimensión mejora, qué dimensión retrocede, si el retroceso es aceptable y quién lo firma». Así, la human evaluation deja de ser solo una exhibición al final del artículo y se convierte en un requisito obligatorio para la liberación. Por ello, la rúbrica de la clase ofrece también una plantilla de ingeniería de uso general: primero se separan los modos de falla y después se asigna a cada dimensión un evaluador calificado y una ruta de escalamiento.

\subsection{Del modelo de lenguaje general al alineamiento con el dominio médico}

La diapositiva sobre alineamiento plantea el siguiente paso: el modelo base ya tiene amplias capacidades lingüísticas y de conocimiento; ¿cómo se adapta al formato y a los criterios de evaluación de las preguntas médicas? Aquí, la «domain alignment» no consiste solo en agregar corpus especializado, sino en que el modelo aprenda, ante instrucciones médicas, a seleccionar el conocimiento pertinente, organizar respuestas largas y satisfacer en la mayor medida posible la rúbrica multidimensional descrita arriba.

\lecturefigure{slide-22-aligning-medical-domain.jpg}{Alineamiento de modelos de lenguaje de gran tamaño con el dominio médico}{Diapositiva V022 recuperada de la grabación oficial; Stanford CS25 Lecture 17}

\teachervoice{En la sesión de preguntas, el ponente reconoció explícitamente que el human evaluation set todavía es pequeño: unas ciento cuarenta preguntas largas. La costosa anotación de expertos es a la vez una limitación de la investigación y un recurso público que vale la pena construir; por ello, las figuras y tablas siguientes deben entenderse como evidencia humana valiosa pero preliminar, no como una garantía estadística que abarque todos los contextos clínicos.}

\begin{warningbox}{La identidad del evaluador determina qué preguntas puede responder}
Un paciente puede juzgar si una respuesta es comprensible y útil, pero no necesariamente si los detalles concuerdan con el consenso clínico; un médico puede juzgar el contenido médico, pero no necesariamente representa las necesidades de comunicación de todos los pacientes. La doble evaluación no consiste en que una sustituya a la otra, sino en colocar cada epistemic role en el lugar que le corresponde.
\end{warningbox}

\subsection{Resumen de la sección}

La evaluación de respuestas largas debe dividirse en varias dimensiones de riesgo y de utilidad, y debe distinguir las responsabilidades de juicio de los médicos y de los usuarios comunes. Este marco proporciona el sistema de coordenadas para los resultados de Med-PaLM que siguen, y muestra también que el cuello de botella de la investigación sobre modelos clínicos no son solo los datos de entrenamiento, sino también los datos de evaluación humana de alta calidad y verificables.
\section{PaLM, Flan-PaLM y Med-PaLM}

Una vez establecido el marco de evaluación, la clase compara tres niveles de modelos: PaLM aporta capacidades básicas a gran escala, Flan-PaLM mejora el seguimiento de instrucciones mediante instruction tuning multitarea y Med-PaLM realiza la alineación con el dominio mediante exemplars médicos e instruction prompt tuning. Cada paso modifica un objeto distinto; entender esta cadena es más importante que memorizar una cifra de una tabla de clasificación.

\subsection{División del trabajo entre la escala y el instruction tuning}

La diapositiva de PaLM presenta la escala y el contexto arquitectónico del Pathways Language Model; la diapositiva de resultados muestra que ampliar el modelo produce ganancias notables en diversas tareas. La diapositiva de Flan-PaLM, por su parte, destaca el instruction fine-tuning: el modelo se entrena con un gran número de tareas descritas en lenguaje natural y aprende a «entender la tarea y responder según lo solicitado», en lugar de solo seguir prediciendo el siguiente token.

\lecturefigure{slide-23-palm-architecture-scale.jpg}{PaLM: modelo de lenguaje base a gran escala}{Diapositiva recuperada de la grabación oficial V023; Stanford CS25 Lecture 17}

\lecturefigure{slide-24-palm-benchmark-results.jpg}{Ganancias por escala de PaLM en benchmarks generales}{Diapositiva recuperada de la grabación oficial V024; Stanford CS25 Lecture 17}

\lecturefigure{slide-25-flan-palm.jpg}{Flan-PaLM: aprender a seguir tareas mediante instruction tuning}{Diapositiva recuperada de la grabación oficial V025; Stanford CS25 Lecture 17}

\begin{knowledgebox}{Términos clave: tres tipos de «ajuste»}
\term{continued pretraining} sigue optimizando los parámetros del modelo para que absorba un corpus del dominio; \term{instruction tuning} actualiza los parámetros del modelo con muchas tareas y ejemplos explícitos para que siga instrucciones en lenguaje natural; \term{instruction prompt tuning} congela el modelo grande, aprende únicamente un breve soft prompt continuo y lo combina con exemplars médicos para guiar al modelo. Los tres difieren en costo, controlabilidad y evidencia.
\end{knowledgebox}

\subsection{Resultados en preguntas médicas de opción múltiple y curvas de Scaling}

Las tres diapositivas de benchmarks cubren, respectivamente, MedQA, MedMCQA/PubMedQA y los temas clínicos de MMLU. Primero conviene observar la tendencia relativa: el Flan-PaLM más grande supera con claridad a las líneas base especializadas anteriores, lo que indica que la escala y el instruction tuning de un modelo general se transfieren a las preguntas médicas de opción múltiple. Después hay que ver las limitaciones: las preguntas suelen ser cerradas, las respuestas son breves y la exactitud no verifica si la explicación que da el modelo es segura.

\lecturefigure{slide-26-medqa-sota.jpg}{Resultados state-of-the-art en MedQA (estilo USMLE)}{Diapositiva recuperada de la grabación oficial V026; Stanford CS25 Lecture 17}

\lecturefigure{slide-27-medmcqa-pubmedqa-results.jpg}{Comparación de resultados en MedMCQA y PubMedQA}{Diapositiva recuperada de la grabación oficial V027; Stanford CS25 Lecture 17}

\lecturefigure{slide-28-mmlu-clinical-results.jpg}{Resultados de los modelos en los temas clínicos de MMLU}{Diapositiva recuperada de la grabación oficial V028; Stanford CS25 Lecture 17}

La diapositiva de Scaling plots coloca varios tamaños de modelo en el eje horizontal; el eje vertical es la exactitud en distintas tareas médicas. La tendencia principal es que el desempeño aumenta con la escala, pero la pendiente y el punto de saturación varían según la tarea; esto indica que la scale es una variable fuerte, pero no la única. Si dos modelos tienen datos de instruction tuning distintos, la diferencia entre sus curvas tampoco puede atribuirse por completo al número de parámetros.

\lecturefigure{slide-29-scaling-plots.jpg}{Tendencias de Scaling en la respuesta a preguntas médicas}{Diapositiva recuperada de la grabación oficial V029; Stanford CS25 Lecture 17}

Al comparar curvas de scaling también hay que controlar los presupuestos de entrenamiento y de inferencia. Un modelo más grande puede haber visto más tokens, usar una instruction mixture más sólida o recibir más ejemplos en el prompt few-shot; si estas condiciones difieren, el «número de parámetros» del eje horizontal no es lo único que cambia. El contexto médico añade además restricciones de latencia, costo y privacidad: un modelo enorme en la nube que obtiene unos cuantos puntos porcentuales más en un benchmark no necesariamente supera a un modelo pequeño que puede ejecutarse en un entorno controlado del hospital, citar las guías internas y abstenerse de responder de manera estable. La pregunta planteada en clase, «¿todavía se necesitan clinical LM?», recuerda precisamente que la Pareto frontier incluye al menos la exactitud, el riesgo de las respuestas largas, el costo de inferencia, la gobernanza de datos y la facilidad de actualización, y no una sola tabla de clasificación.

\begin{warningbox}{Lo que no puede inferirse al leer la figura}
Una scaling curve ascendente no demuestra que el modelo haya desarrollado un razonamiento clínico confiable, ni que seguir aumentando los parámetros elimine necesariamente los daños y los sesgos. Solo indica que, con los datos, los prompts y las reglas de calificación actuales, la escala está correlacionada con la exactitud en la tarea; la evidencia sobre las respuestas largas debe medirse por separado.
\end{warningbox}

\subsection{Selective prediction: permitir que el modelo se abstenga}

La selective prediction (predicción selectiva) hace que el modelo responda solo las muestras de alta confianza. Sea $D$ el conjunto de muestras, $\hat y_i$ la respuesta del modelo, $s_i$ su puntuación de confianza y $\tau$ el umbral. La cobertura y la exactitud condicional son, respectivamente,
$$
\operatorname{coverage}(\tau)=\frac{1}{|D|}\sum_{i\in D}\mathbf{1}[s_i\ge\tau],
\qquad
\operatorname{acc}(\tau)=
\frac{\sum_i\mathbf{1}[s_i\ge\tau]\mathbf{1}[\hat y_i=y_i]}
{\sum_i\mathbf{1}[s_i\ge\tau]}.
$$
Aquí $D$ representa las muestras de evaluación, $s_i$ la confianza de la $i$-ésima respuesta, $\tau$ el umbral de abstención, $\mathbf{1}[\cdot]$ la función indicadora y $y_i$ la respuesta correcta. Elevar $\tau$ suele reducir la coverage y aumentar la exactitud de las respuestas que se conservan, pero esto supone que la confianza está correlacionada con los errores reales.

\lecturefigure{slide-30-selective-prediction.jpg}{Selective prediction: compromiso entre cobertura y exactitud}{Diapositiva recuperada de la grabación oficial V030; Stanford CS25 Lecture 17}

\begin{lstlisting}[caption={Flujo de umbralización de la predicción selectiva (esquemático)}]
def selective_answer(question, threshold):
    answer, confidence = model.generate_with_confidence(question)
    if confidence < threshold:
        return "ABSTAIN"
    return answer
\end{lstlisting}

\begin{importantbox}{Condiciones para que la abstención sea útil}
La abstention solo tiene valor si la calibración de la confianza es efectiva, si los casos en que el modelo se abstiene pueden pasar a una persona o a un sistema más confiable y si el servicio puede aceptar una cobertura baja. Un sistema médico también debe revisar los «errores en los que el modelo tiene más confianza», porque una mala calibración conserva justamente las muestras de riesgo.
\end{importantbox}

La predicción selectiva también requiere definir qué ocurre después de la abstención. Si una respuesta de baja confianza simplemente se devuelve en blanco, el usuario puede recurrir a fuentes de información menos confiables; si el sistema escala automáticamente el caso a un clinician, hay que medir la cola de casos escalados, el tiempo de respuesta y la carga de trabajo. El umbral tampoco puede elegirse solo con un benchmark estático, porque la distribución real de las preguntas se desplaza, y los fármacos nuevos, las guías nuevas y los casos raros pueden producir un exceso de confianza sistemático. Un ciclo de monitoreo completo debe registrar la coverage, la tasa de error, el error de calibración, los motivos de abstención y el resultado de la intervención humana, y revisarlos estratificados por servicio clínico o por población. Solo así se sabe si «responder menos» realmente transfiere el riesgo a un proceso más confiable, en lugar de excluir del servicio a los usuarios con casos difíciles.

\subsection{Por qué sigue siendo necesaria la alineación con el dominio médico}

La diapositiva de Human evaluation limitations reconoce que las respuestas abiertas son difíciles de evaluar a escala; la siguiente explica por qué no se puede uno detener en Flan-PaLM: aunque el modelo general es muy bueno en preguntas de opción múltiple, sus respuestas largas todavía pueden apartarse del consenso médico, omitir información esencial o no ser adecuadas para que las lea un paciente. Los datos médicos y las rúbricas de expertos no buscan volver a entrenar todo el conocimiento del mundo, sino cambiar la conducta al responder.

\lecturefigure{slide-31-human-eval-limitations.jpg}{Limitaciones de escala y costo en la evaluación humana de respuestas largas}{Diapositiva recuperada de la grabación oficial V031; Stanford CS25 Lecture 17}

\lecturefigure{slide-32-need-domain-data.jpg}{Por qué siguen siendo necesarios los datos del dominio médico y las señales de alineación}{Diapositiva recuperada de la grabación oficial V032; Stanford CS25 Lecture 17}

\subsection{Instruction prompt tuning y el flujo de Med-PaLM}

La diapositiva de Prompt tuning muestra que los parámetros del modelo grande permanecen congelados y que lo que se entrena es un conjunto de prompt vectors continuos; los exemplars médicos, por su parte, ofrecen a la tarea una demostración de formato y de semántica. Al combinar ambos se obtiene el flujo de alineación de Med-PaLM. No se trata del RLHF que describía el borrador anterior de la clase, y en la clase tampoco se introdujeron fairness prompts bilingües ni una retroalimentación semanal de clinician review; una reconstrucción confiable debe conservar solo los pasos que respaldan las diapositivas y el artículo.

\lecturefigure{slide-33-prompt-tuning.jpg}{Instruction prompt tuning: congelar el modelo y aprender un prompt suave}{Diapositiva recuperada de la grabación oficial V033; Stanford CS25 Lecture 17}

\lecturefigure{slide-34-medpalm-pipeline.jpg}{Combinación de exemplars médicos y prompt tuning para obtener Med-PaLM}{Diapositiva recuperada de la grabación oficial V034; Stanford CS25 Lecture 17}

Sea $f_{\theta}$ el modelo de lenguaje congelado, $P\in\mathbb{R}^{m\times d}$ el prompt suave aprendible, $X$ el embedding de los tokens de entrada y $Y$ la respuesta objetivo. El objetivo de entrenamiento puede escribirse como
$$
P^{*}=\arg\min_{P}\; -\sum_{t=1}^{|Y|}\log p_{\theta}(Y_t\mid P,X,Y_{<t}).
$$
Aquí $\theta$ son los parámetros congelados, $m$ es la longitud del prompt suave, $d$ es la dimensión del embedding y $Y_t$ es el $t$-ésimo token objetivo. Las actualizaciones se concentran en $P$, por lo que la eficiencia en parámetros es alta; sin embargo, lo que aprende el prompt suave sigue dependiendo de los exemplars, de las tareas de entrenamiento y de los criterios de evaluación.

\teachervoice{El ponente llama a esta parte «putting it all together»: el modelo base aporta un conocimiento amplio, el instruction tuning aporta el seguimiento de tareas, y el prompt tuning médico junto con los exemplars aportan la conducta propia del dominio. En la clase no se presentó ninguno de estos pasos como garantía de seguridad; más adelante todavía hay que compararlos eje por eje con las clinician answers.}

\begin{warningbox}{No describir el soft prompt como un prompt de sistema legible}
Los prompt vectors continuos son parámetros aprendibles en el espacio de embedding del modelo y no necesariamente corresponden a palabras legibles para una persona. Son distintos de un prompt discreto como añadir «responde con cautela» en la ventana de chat, y tampoco ofrecen de forma automática reglas médicas interpretables.
\end{warningbox}

Desde el punto de vista de la reproducibilidad, este pipeline requiere registrar al menos cuatro tipos de configuración: el checkpoint del modelo base, la versión del instruction tuning, el conjunto de exemplars médicos y los parámetros del soft prompt. Si solo se publica el nombre del modelo final, no puede determinarse si la mejora proviene de la escala, de la mezcla de tareas o de la selección de ejemplos. En particular, en la respuesta a preguntas médicas, los exemplars few-shot pueden filtrar el formato de la respuesta o una distribución implícita de etiquetas; por ello, el umbral debe volver a elegirse en un conjunto de preguntas independiente y debe reportarse la prompt sensitivity. Que el método sea eficiente en parámetros no significa que la evaluación sea sencilla: aunque solo se actualice una pequeña cantidad de vectores, la conducta del modelo todavía puede cambiar de forma notable según la redacción, la longitud de la pregunta y el contexto del usuario. Las notas de la clase separan estos componentes para que el lector pueda localizar las fallas, en lugar de atribuir toda la conducta a la etiqueta única «Med-PaLM».

\subsection{Resumen de la sección}

PaLM, Flan-PaLM y Med-PaLM aportan, respectivamente, la escala base, el seguimiento general de instrucciones y la conducta propia del dominio médico. Los resultados en preguntas de opción múltiple y de scaling demuestran que los modelos grandes tienen una fuerte capacidad de transferencia; la selective prediction ofrece un mecanismo de abstención, y el instruction prompt tuning logra la alineación con el dominio con relativamente pocos parámetros entrenables. Aun así, estos mecanismos deben ponerse a prueba mediante la evaluación humana de respuestas largas.
\section{Evaluación humana multidimensional de Med-PaLM}

Una vez armada la cadena de modelos, la clase vuelve a la pregunta más importante sobre la evidencia: ¿Med-PaLM solo es más fuerte en preguntas de opción múltiple, o también ha reducido la distancia con los médicos en el rigor científico, el riesgo y la utilidad de las respuestas largas reales? Las nueve diapositivas de resultados que siguen no deben leerse solo por la altura de las barras o los colores, sino eje por eje.

\subsection{Consenso científico y evidencia de razonamiento}

La diapositiva de Scientific consensus compara Flan-PaLM, Med-PaLM y las clinician answers. La tendencia general es que Med-PaLM mejora de forma notable y se acerca a los expertos, aunque todavía queda una brecha. La siguiente diapositiva separa comprehension, knowledge retrieval y reasoning; en la sesión de preguntas se señaló en particular que la evidencia correcta y la evidencia incorrecta no son complementarias en el sentido de $1-x$: una misma respuesta puede contener al mismo tiempo hechos correctos e inferencias erróneas.

\lecturefigure{slide-35-clinical-consensus.jpg}{Concordancia con el consenso científico y clínico}{Diapositiva V035 recuperada de la grabación oficial; Stanford CS25 Lecture 17}

\lecturefigure{slide-36-comprehension-retrieval-reasoning.jpg}{Evidencia de comprensión, recuperación de conocimiento médico y razonamiento}{Diapositiva V036 recuperada de la grabación oficial; Stanford CS25 Lecture 17}

\teachervoice{Un estudiante preguntó en clase si los dos indicadores de comprehension eran complementarios, y el ponente respondió claramente «no». La comprensión correcta y la incorrecta pueden aparecer a la vez; por ejemplo, una respuesta puede identificar primero la enfermedad de forma correcta y luego añadir una inferencia irrelevante en la recomendación de tratamiento. Es una corrección muy importante para leer gráficas complejas de human-evaluation.}

\begin{knowledgebox}{Lectura de la figura: por qué los indicadores pueden ser verdaderos a la vez}
Una respuesta larga se compone de muchas proposition. Si una parte de ellas es correcta y otra parte es incorrecta, tanto «hay evidencia de comprensión correcta» como «hay evidencia de comprensión incorrecta» pueden ser verdaderas. Reducir todo el texto a correcto o incorrecto con un solo accuracy pierde este estado mixto; el valor de la evaluación eje por eje está precisamente en exponer las fallas locales.
\end{knowledgebox}

\subsection{Las respuestas más largas también pueden producir más errores}

La diapositiva de Incorrect or missing content revela un efecto secundario de la alineación: Med-PaLM se orienta a generar respuestas más completas y detalladas, pero el contenido adicional también amplía la oportunidad de que aparezcan afirmaciones irrelevantes o erróneas. El ponente explica que, en algunos ejes, el hecho de que Flan-PaLM parezca un poco mejor no significa que sea más seguro en conjunto; puede deberse simplemente a que dice menos. La evaluación debe considerar al mismo tiempo la completeness y la precision.

\lecturefigure{slide-37-incorrect-missing-context.jpg}{El doble riesgo del contenido erróneo y del contenido faltante}{Diapositiva V037 recuperada de la grabación oficial; Stanford CS25 Lecture 17}

\begin{importantbox}{La paradoja de la longitud en la generación clínica}
Una respuesta breve puede omitir advertencias críticas; una respuesta larga puede introducir alucinaciones adicionales. El objetivo no es maximizar el número de palabras, sino cubrir la información necesaria para la decisión mientras se controlan las proposition sin sustento. En términos de ingeniería, la «longitud de la respuesta» debe tratarse como una variable de riesgo, no como una señal natural de calidad.
\end{importantbox}

\subsection{Daño y sesgo: las consecuencias de los errores no son equivalentes}

La diapositiva de Harm considera a la vez la probabilidad y la gravedad. En la sesión de preguntas se dio una explicación concreta: el daño grave no proviene solo de recomendar un tratamiento peligroso, sino también de un misdiagnosis, de pasar por alto una condición que pone en riesgo la vida o de no comunicar la gravedad de la enfermedad. La diapositiva de Bias examina si la respuesta introduce sesgos demográficos inapropiados. Ambos aspectos requieren el juicio de expertos del dominio que tomen en cuenta el contexto; no pueden sustituirse con una lista negra de palabras clave.

\lecturefigure{slide-38-extent-likelihood-harm.jpg}{Grado y probabilidad del daño potencial}{Diapositiva V038 recuperada de la grabación oficial; Stanford CS25 Lecture 17}

\lecturefigure{slide-39-potential-bias.jpg}{Evaluación del sesgo potencial en las respuestas}{Diapositiva V039 recuperada de la grabación oficial; Stanford CS25 Lecture 17}

\begin{warningbox}{El daño no consiste en «si la respuesta contiene palabras peligrosas»}
Aunque el modelo no dé directamente una recomendación peligrosa, puede provocar un retraso en el tratamiento por un diagnóstico omitido, por subestimar la gravedad o por tranquilizar indebidamente. A la inversa, mencionar el cáncer o la muerte no constituye daño de forma automática; lo determinante es si el contenido corresponde a la pregunta y a la evidencia, y cómo podría cambiar la conducta del usuario.
\end{warningbox}

\teachervoice{El ponente concretó los ejemplos de harm grave como no identificar una patología mortal, por ejemplo un cancer, o no comunicar adecuadamente la gravedad. También señaló que el tema tiene muchos matices y que el artículo se apoya en el marco de riesgo de la AHRQ; la gráfica de la clase es solo una visión general de alto nivel y estas notas no deben reescribirla como una norma completa de seguridad clínica.}

\subsection{Lo que ve el usuario común es la helpfulness y la comprensibilidad}

La diapositiva de resultados de Layperson se centra en la intención y en el grado de ayuda. La tendencia aproximada que se presentó en clase es la siguiente: alrededor del 60\,\% de las respuestas de Flan-PaLM se consideran helpful, Med-PaLM sube a cerca del 80\,\% y las clinician answers se acercan al 90\,\%. Lo importante no es tomar estas proporciones como umbrales de despliegue, sino ver que la alineación médica puede reducir la brecha de comunicación, aunque sin alcanzar el nivel de los médicos.

\lecturefigure{slide-40-helpfulness-utility.jpg}{Evaluación de usuarios comunes: cobertura de la intención, grado de ayuda y utilidad}{Diapositiva V040 recuperada de la grabación oficial; Stanford CS25 Lecture 17}

\lecturefigure{slide-41-clinical-qualitative-example.jpg}{Caso cualitativo: diferencias entre las respuestas del médico, de Flan-PaLM y de Med-PaLM}{Diapositiva V041 recuperada de la grabación oficial; Stanford CS25 Lecture 17}

La diapositiva del ejemplo cualitativo debe leerse a partir de la forma en que se organiza la información: las respuestas del médico suelen ser precisas y breves, pero la terminología especializada puede dificultar que el paciente las entienda; el modelo puede desarrollar la explicación y añadir contexto, aunque al extenderse también puede introducir más errores. Lo que ambos se complementan es la conversión del lenguaje y la organización de la información, no que el modelo sustituya el juicio médico sin supervisión.

\begin{knowledgebox}{Lectura de la figura: cuatro aspectos que revisar primero al comparar respuestas}
Primero, verifique si la respuesta atiende directamente la intención del usuario; después, si están completos los hechos médicos esenciales; luego, si los detalles añadidos tienen evidencia, y por último, si se indica con claridad el límite a partir del cual hay que acudir a un profesional. La fluidez solo puede considerarse después de estos cuatro aspectos; de lo contrario, la respuesta mejor redactada podría juzgarse erróneamente como la más segura.
\end{knowledgebox}

El caso cualitativo también muestra los puntos ciegos de las metric automáticas. Dos respuestas pueden compartir gran parte del vocabulario y, sin embargo, tener consecuencias completamente distintas por una negación crucial, una dosis o una recomendación de acudir al médico; a la inversa, un médico puede dar la conclusión correcta con una sola abreviatura especializada y la similitud lingüística resultar muy baja. Una evaluación adecuada debe descomponer primero la respuesta en proposition clínicas, comprobar si cada una está respaldada por la pregunta y por el consenso, y después evaluar la organización, el tono y la posibilidad de actuar a partir de ella. Si el sistema se usa para reescribir las notas del médico para el paciente, también debe exigirse que conserve un enlace al texto original y un punto de confirmación humana, de modo que la «versión fácil de entender» no pueda alterar en silencio los hechos médicos. Así, la ventaja del modelo de lenguaje se aprovecha para explicar y convertir, mientras la fuente autorizada sigue siendo rastreable.

\subsection{¿Siguen haciendo falta los modelos de lenguaje clínicos?}

La penúltima diapositiva plantea una pregunta abierta: si los grandes modelos de propósito general ya son tan capaces, ¿siguen haciendo falta modelos de lenguaje especializados en lo clínico (clinical LM)? La clase no dio una respuesta negativa simple. Los modelos pequeños de dominio pueden aportar valor en costo, privacidad, despliegue local o distribuciones de datos específicas; los grandes modelos de propósito general, en cambio, aprovechan un conocimiento más amplio y el scaling. La verdadera pregunta es con qué datos, tareas y restricciones el beneficio marginal del preentrenamiento de dominio supera al de un modelo general más grande.

\lecturefigure{slide-42-clinical-lm-question.jpg}{Pregunta abierta: ¿siguen haciendo falta modelos de lenguaje clínicos especializados?}{Diapositiva V042 recuperada de la grabación oficial; Stanford CS25 Lecture 17}

\lecturefigure{slide-43-clinical-key-takeaways.jpg}{Conclusiones clave y limitaciones de la sección clínica}{Diapositiva V043 recuperada de la grabación oficial; Stanford CS25 Lecture 17}

\teachervoice{El papel que el ponente asigna a las aplicaciones clínicas a corto plazo no es el de un autonomous physician, sino el de augment physicians: convertir el doctor speak difícil de entender en explicaciones más accesibles para el paciente y ayudar a la comunicación entre médico y paciente, entre médicos y entre médicos e investigadores. Al mismo tiempo, dijo que esto todavía «no es perfecto» y que la oportunidad actual está en la complementarity.}

\begin{importantbox}{Conclusiones de la evidencia en la sección clínica}
La evidencia de la clase respalda tres puntos: los grandes modelos de propósito general transfieren con fuerza a las preguntas médicas de opción múltiple; el prompt tuning médico puede mejorar varios indicadores humanos de las respuestas largas, y el valor a corto plazo se parece más al apoyo a la comunicación y al conocimiento. No respalda que «aprobar un examen equivalga a obtener la habilitación clínica», ni que «una respuesta larga, por ser más detallada, sea más segura».
\end{importantbox}

\subsection{Resumen de la sección}

Med-PaLM supera claramente a Flan-PaLM en varias evaluaciones humanas, como el consenso científico, el razonamiento y el grado de ayuda, y reduce la brecha con las respuestas de los médicos; al mismo tiempo, las respuestas detalladas pueden introducir más errores, y el daño y el sesgo siguen requiriendo el juicio de expertos. La conclusión más confiable de la clase es el aumento clínico supervisado, no la sustitución automática.
\section{Secuencias largas de proteínas: Performer y Protein LM}

En la sección clínica, los retos principales eran el conocimiento, la evaluación y el riesgo; al pasar a las proteínas, el primer obstáculo se vuelve la longitud de la secuencia y la complejidad computacional. Una secuencia de aminoácidos puede darse como entrada al modelo igual que una secuencia de tokens, pero las proteínas largas hacen que la matriz de correlaciones $n\times n$ de la attention estándar consuma rápidamente el cómputo y la memoria. Aquí, Performer no es un recurso abstracto de eficiencia, sino la condición previa para poder entrenar con secuencias biológicas más largas.

\subsection{Por qué las proteínas necesitan contexto largo}

La diapositiva de transición de sección cambia la perspectiva del lenguaje de los médicos al lenguaje de las moléculas. La siguiente diapositiva señala que la longitud de las proteínas varía en un intervalo muy amplio, y que aminoácidos lejanos en la secuencia pueden quedar próximos entre sí después del plegamiento y determinar juntos el sitio activo. Si solo se observa una ventana pequeña y fija, se puede perder esta dependencia; si se usa directamente la self-attention global, hay que pagar una complejidad cuadrática.

\lecturefigure{slide-44-section-proteins.jpg}{Cambio de sección: aplicaciones en proteínas}{Diapositiva V044 recuperada de la grabación oficial; Stanford CS25 Lecture 17}

\lecturefigure{slide-45-long-biological-sequences.jpg}{Problemas de modelado y de cómputo que plantean las secuencias biológicas largas}{Diapositiva V045 recuperada de la grabación oficial; Stanford CS25 Lecture 17}

La scaled dot-product attention estándar es
$$
\operatorname{Attn}(Q,K,V)=
\operatorname{softmax}\!\left(\frac{QK^{\top}}{\sqrt d}\right)V.
$$
Donde $Q,K,V\in\mathbb{R}^{n\times d}$ son, respectivamente, query, key y value; $n$ es la longitud de la secuencia y $d$ es la dimensión de cada head. El cuello de botella proviene de $QK^{\top}\in\mathbb{R}^{n\times n}$: aunque al final solo se necesite una salida de $n\times d$, hay que procesar, de forma explícita o implícita, todos los pares de posiciones.

\begin{warningbox}{La dificultad de las secuencias largas no es solo «falta de memoria de GPU»}
La attention cuadrática aumenta a la vez el tiempo de entrenamiento, las activaciones intermedias y la presión de comunicación. El truncamiento simple ahorra cómputo, pero puede eliminar motifs lejanos o relaciones entre dominios estructurales; la attention por bloques, a su vez, puede separar posiciones que en realidad están relacionadas. El objetivo de los métodos de eficiencia es lograr un compromiso controlado entre el presupuesto de cómputo y la información global.
\end{warningbox}

\subsection{Performer: reordenar el cálculo con kernel features}

La subsección anterior ya ubicó el cuello de botella en la score matrix completa de $n\times n$; esta sección responde, además, cómo cambiar el orden del cálculo sin fragmentar la secuencia. La diapositiva de Performer presenta la idea de FAVOR+: aproximar el kernel softmax como un producto interno de características, de modo que las dependencias globales todavía puedan propagarse a través de un estadístico comprimido. Al leer la figura hay que observar al mismo tiempo el reordenamiento matemático, el error de la aproximación aleatoria y las curvas en hardware real, en lugar de tomar «linear attention» solo como un lema de complejidad.
$$
\exp(q^{\top}k)\approx \phi(q)^{\top}\phi(k),
$$
y después, por asociatividad, se agrega primero $\phi(K)^{\top}V$, en lugar de construir la attention matrix completa. Si se omiten los detalles de normalización, la salida de una sola query puede escribirse como
$$
\operatorname{Attn}(q_i,K,V)\approx
\frac{\phi(q_i)^{\top}\left(\sum_j \phi(k_j)v_j^{\top}\right)}
{\phi(q_i)^{\top}\left(\sum_j \phi(k_j)\right)}.
$$
Donde $\phi(\cdot)$ es un mapeo aleatorio de características positivas, y $q_i,k_j,v_j$ representan, respectivamente, la $i$-ésima query y el $j$-ésimo key/value. Después del reordenamiento ya no se guardan $n^2$ pairwise scores, y el tiempo y el espacio pueden crecer de forma aproximadamente lineal con $n$; el costo es la varianza que introduce la kernel approximation y la complejidad de implementación.

\lecturefigure{slide-46-performer-kernel-attention.jpg}{Performer: aproximación de la softmax attention con características aleatorias}{Diapositiva V046 recuperada de la grabación oficial; Stanford CS25 Lecture 17}

\lecturefigure{slide-47-performer-speed-complexity.jpg}{Comparación de velocidad, memoria y complejidad de Performer}{Diapositiva V047 recuperada de la grabación oficial; Stanford CS25 Lecture 17}

\begin{lstlisting}[caption={Orden de cálculo de la attention linealizada (esquema)}]
def performer_attention(query, key, value, phi):
    q_features = phi(query)
    k_features = phi(key)
    kv_summary = k_features.T @ value
    k_summary = k_features.sum(axis=0)
    numerator = q_features @ kv_summary
    denominator = q_features @ k_summary
    return numerator / denominator[:, None]
\end{lstlisting}

\begin{knowledgebox}{Lectura de la figura: cómo leer la gráfica de complejidad}
Primero se fijan la dimensión del modelo y el lote, y luego se aumenta la sequence length a lo largo del eje horizontal; las curvas de tiempo y de memoria de la attention estándar suben más rápido, mientras que el crecimiento de Performer se acerca más a lo lineal. La figura respalda que «las secuencias largas son más escalables», pero no demuestra que sea más rápido para todas las longitudes, todo el hardware y todos los objetivos de precisión; el número de características aleatorias, la implementación del kernel y los términos constantes modifican el punto de cruce.
\end{knowledgebox}

Desde el punto de vista de la implementación, la ganancia de la attention lineal depende de que los tensores intermedios se calculen realmente en el orden reordenado. Si, para depurar, el código todavía genera explícitamente los pairwise scores, la complejidad teórica no se materializa; si el número de features aleatorias es demasiado grande, el término constante de $n r d$ también puede superar a una attention estándar bien optimizada. En cambio, en una tarea como la de proteínas, con una distribución de longitudes muy amplia, puede elegirse la implementación por bucket de longitud: usar attention exacta para secuencias cortas, cambiar a Performer para las largas y comparar la salida aproximada con la exacta en un conjunto de muestras que el presupuesto permita. Esta validación vincula «más rápido» con «sesgo aceptable» y evita medir solo el rendimiento ignorando los cambios en el comportamiento del modelo.

\teachervoice{El ponente coloca Performer al inicio de la sección de proteínas porque las secuencias biológicas chocan con el costo de la attention cuadrática «antes de empezar a hacer razonamiento biológico». Lo importante en la clase no es memorizar el nombre FAVOR+, sino ver que la distribución de longitudes de los datos del dominio, a su vez, impulsa nuevos attention algorithms.}

\subsection{Qué aprende un Protein LM}

La diapositiva de resultados muestra la ganancia de los modelos de lenguaje en protein sequence modeling, y la diapositiva de attention visualization dibuja la atención entre distintos amino acids. Una interpretación razonable es la siguiente: el preentrenamiento obliga al modelo a aprovechar motifs conservados, el entorno químico local y las relaciones entre residuos más lejanos para predecir las posiciones enmascaradas, por lo que en las representaciones aparecen patrones relacionados con la estructura o la función.

\lecturefigure{slide-48-protein-lm-results.jpg}{Resultados experimentales de protein language modeling}{Diapositiva V048 recuperada de la grabación oficial; Stanford CS25 Lecture 17}

\lecturefigure{slide-49-amino-acid-attention.jpg}{Attention pattern entre aminoácidos y dependencias complejas}{Diapositiva V049 recuperada de la grabación oficial; Stanford CS25 Lecture 17}

\begin{warningbox}{Un attention map no es una explicación causal}
Un attention weight alto indica que ciertas posiciones interactúan con fuerza en el cálculo del modelo; no equivale a que ese par de residuos esté en contacto directo en la proteína real, ni demuestra que provoquen alguna función. Una conclusión biológica confiable requiere además datos estructurales, mutagenesis u otros experimentos. La visualización sirve sobre todo para plantear hipótesis y revisar el modelo, no para sustituir la validación experimental.
\end{warningbox}

\begin{importantbox}{El sentido del caso Performer en el curso}
Cuando los datos son largos por naturaleza, la arquitectura del modelo y la complejidad del algoritmo determinan qué preguntas científicas pueden plantearse. Performer amplía el contexto procesable mediante una kernel attention aproximada; el protein LM, por su parte, aprovecha ese contexto para aprender regularidades estadísticas. El avance en eficiencia y el descubrimiento en el dominio son dos niveles de evidencia y no deben fusionarse en una sola conclusión.
\end{importantbox}

\subsection{Resumen de la sección}

Las secuencias de proteínas exponen de manera muy directa el costo de contexto largo del Transformer. Performer evita la attention matrix completa de $n\times n$ mediante características aleatorias y la asociatividad, lo que hace más viable el entrenamiento con secuencias largas; las representaciones y los attention patterns del protein LM pueden aportar indicios biológicos, pero todavía requieren validación estructural y experimental.
\section{ProtNLM: generación de descripciones funcionales en lenguaje natural a partir de secuencias de aminoácidos}

La sección anterior seguía centrada en la representación de secuencias y en la masked prediction; ProtNLM lleva un paso más allá la salida y la convierte en lenguaje natural. El valor de esta tarea radica en que la protein function en las bases de datos ya existe como texto libre; si el modelo puede generar a partir de la sequence descripciones legibles y que se puedan buscar, puede ayudar a los investigadores a proponer funciones candidatas para la gran cantidad de proteínas que no están suficientemente anotadas.

\subsection{Por qué la anotación de proteínas es una tarea de lenguaje}

La diapositiva de título presenta ProtNLM, y las dos siguientes explican la necesidad: las secuencias de proteínas producidas por la secuenciación crecen con rapidez, mientras que la anotación experimental manual es costosa y lenta; bases de datos como UniProt contienen una gran cantidad de descripciones en texto libre, pero su profundidad y confiabilidad son desiguales. Lo que el modelo debe aprender no es a «ponerle un título vistoso a una secuencia», sino a asignar la sequence evidence a enunciados en lenguaje sobre función, familia, localización y procesos biológicos.

\lecturefigure{slide-50-protnlm-title.jpg}{ProtNLM: Model-based Natural Language Protein Annotation}{Diapositiva V050 recuperada de la grabación oficial; Stanford CS25 Lecture 17}

\lecturefigure{slide-51-protein-function-demand.jpg}{Enorme cantidad de secuencias de proteínas y necesidad de anotación funcional}{Diapositiva V051 recuperada de la grabación oficial; Stanford CS25 Lecture 17}

\lecturefigure{slide-52-uniprot-free-text.jpg}{Anotaciones de proteínas en texto libre en UniProt}{Diapositiva V052 recuperada de la grabación oficial; Stanford CS25 Lecture 17}

\begin{knowledgebox}{Glosario: niveles de evidencia de la protein annotation}
\term{sequence annotation} predice etiquetas o descripciones a partir de la secuencia de aminoácidos; \term{curated annotation} la elaboran expertos combinando experimentos y literatura; \term{automatic annotation} se propaga en bloque con base en similitud o en modelos; \term{experimental validation} confirma la función mediante experimentos bioquímicos, estructurales o celulares. ProtNLM mejora sobre todo la interfaz de información entre los dos primeros niveles; no puede elevar el texto generado a la categoría de hecho experimental.
\end{knowledgebox}

\subsection{Sequence-to-language: conectar dos vocabularios}

La diapositiva de sequence-to-language coloca los amino acid token de entrada y los natural-language token de salida en un mismo marco encoder--decoder; el ejemplo de captioning muestra que el modelo puede generar descripciones funcionales coherentes. Lo esencial aquí no es la fluidez del lenguaje, sino si la descripción contiene información de family, domain, localization o process congruente con la secuencia de entrada.

\lecturefigure{slide-53-sequence-to-language.jpg}{Tarea encoder--decoder de secuencia de proteína a lenguaje natural}{Diapositiva V053 recuperada de la grabación oficial; Stanford CS25 Lecture 17}

\lecturefigure{slide-54-protein-captioning.jpg}{Ejemplo de entrada y descripción generada en el captioning de proteínas}{Diapositiva V054 recuperada de la grabación oficial; Stanford CS25 Lecture 17}

Dada una secuencia de aminoácidos $S=(s_1,\ldots,s_n)$ y una descripción objetivo $Y=(y_1,\ldots,y_T)$, el objetivo de generación condicional es
$$
\mathcal{L}_{\text{caption}}=-\sum_{t=1}^{T}\log p(y_t\mid S,y_{<t}).
$$
Donde $s_i$ es el $i$-ésimo amino acid, $y_t$ es el $t$-ésimo token de lenguaje natural, y $n$ y $T$ son las longitudes de la entrada y de la salida, respectivamente. Este objetivo recompensa la generación congruente con el texto de referencia, pero si las propias etiquetas de la base de datos faltan o tienen mucho ruido, el modelo también hereda esa incertidumbre.

\begin{warningbox}{Que la generación suene bien no significa que la anotación sea verdadera}
Un objetivo de texto libre recompensa las formulaciones frecuentes, y el modelo puede apoyarse en el conocimiento previo de la familia para generar descripciones que suenan razonables pero que no se han confirmado experimentalmente. La evaluación debe distinguir lexical similarity, exactitud de las entidades especializadas, congruencia del nivel funcional y predicciones novedosas; toda descripción automática debe conservar su provenance y su incertidumbre.
\end{warningbox}

La partición del conjunto de entrenamiento es especialmente delicada aquí. Si proteínas altamente homólogas aparecen tanto en entrenamiento como en prueba, el modelo puede repetir descripciones de familia memorizando vecinos cercanos, lo que parece una inferencia de la función a partir de la sequence. Una evaluación más estricta agrupa por sequence identity o por protein family, reserva las familias lejanas para el conjunto de prueba y reporta por separado las funciones comunes y las raras. El texto generado también puede descomponerse en campos estructurados, por ejemplo family, catalytic activity, subcellular location y evidence qualifier; la verificación campo por campo se acerca más al uso científico que una sola puntuación global de BLEU o ROUGE. La interfaz en lenguaje natural es cómoda, pero cuanto más cómoda es, más necesario resulta escribir de forma explícita el nivel de evidencia en la salida.

\subsection{Entrenamiento multitarea con T5 y resultados}

La diapositiva de T5 muestra que ProtNLM no se limita a un captioning unidireccional, sino que unifica span corruption, sequence-to-text y tareas relacionadas en un formato text-to-text. El entrenamiento multitarea permite que el modelo aprenda al mismo tiempo los patrones locales de las proteínas y su expresión en lenguaje. La diapositiva de resultados resume las mejoras del modelo en la evaluación interna y en el juicio humano; deben entenderse como evidencia de calidad y cobertura de la anotación, no como confirmación experimental de la función de todas las proteínas nuevas.

\lecturefigure{slide-55-protnlm-t5-tasks.jpg}{Preentrenamiento al estilo T5 y objetivos multitarea de ProtNLM}{Diapositiva V055 recuperada de la grabación oficial; Stanford CS25 Lecture 17}

\lecturefigure{slide-56-protnlm-results.jpg}{Resultados de ProtNLM en la generación de descripciones de proteínas}{Diapositiva V056 recuperada de la grabación oficial; Stanford CS25 Lecture 17}

\begin{knowledgebox}{Lectura de la figura: los resultados deben separarse en «cobertura» y «confiabilidad»}
Un modelo automático puede dar con rapidez descripciones para más secuencias; ese es el beneficio de coverage. La congruencia de las descripciones con las anotaciones de referencia o con el juicio de expertos es el beneficio de quality. Ambos deben considerarse a la vez: ampliar solo la cobertura propaga errores en bloque, y buscar solo una precision altísima puede llevar a rechazar una gran cantidad de proteínas desconocidas.
\end{knowledgebox}

\subsection{Caso CRISPR-Cas9: cómo las descripciones generadas sirven al descubrimiento}

Ya se explicó cómo genera el modelo las descripciones de proteínas; esta sección coloca la interfaz dentro de un proceso de descubrimiento real. La diapositiva de aplicación muestra el espacio de proteínas relacionadas con CRISPR-Cas9; los investigadores no se enfrentan a una respuesta conocida, sino a una gran cantidad de candidatas que pueden tener mecanismos similares, tamaños distintos o características de targeting diferentes. Que el modelo genere u organice anotaciones en lenguaje natural puede ayudar a buscar secuencias, comparar familias y formular hipótesis experimentales. El flujo de trabajo correcto es «el modelo filtra y explica candidatas — el investigador revisa la evidencia de bases de datos y estructural — se diseña el experimento — se valida la función», y no «el modelo escribe una descripción — la función ya está descubierta».

\lecturefigure{slide-57-crispr-cas9-application.jpg}{Ejemplo de aplicación: búsqueda de proteínas relacionadas con CRISPR-Cas9}{Diapositiva V057 recuperada de la grabación oficial; Stanford CS25 Lecture 17}

\teachervoice{El ponente usa el caso de CRISPR para mostrar el valor científico de la interfaz en lenguaje natural: los investigadores no tienen que limitarse a buscar vecinos cercanos en el espacio de embedding, sino que pueden organizar las candidatas mediante descripciones funcionales. Aun así, la formulación en clase sigue siendo annotation y discovery support; no trata las oraciones generadas como wet-lab proof.}

\begin{importantbox}{El lugar que corresponde a ProtNLM}
ProtNLM conecta las secuencias de proteínas con conocimiento legible por personas, lo que hace más eficientes la búsqueda, la síntesis y la generación de candidatas. Se parece sobre todo a una «interfaz de información científica»: la entrada es una sequence y la salida es una hypothesis legible; la verdadera función biológica la siguen determinando la provenance de las bases de datos, la literatura y el ciclo experimental.
\end{importantbox}

\subsection{Resumen de la sección}

ProtNLM extiende el aprendizaje de representaciones de proteínas a sequence-to-language y, mediante entrenamiento multitarea al estilo T5, genera descripciones funcionales en texto libre. Puede ampliar la cobertura de la anotación, mejorar la búsqueda de conocimiento y apoyar el descubrimiento de candidatas, pero la calidad del texto generado y la veracidad de la función biológica deben evaluarse por separado.
\section{DeepConsensus: corrección de lecturas de secuenciación con Transformer}

Al entrar en genomics, la clase aborda primero un problema de «calidad de la entrada». El análisis genómico depende de los read de secuenciación; si un read contiene errores, se ven afectados el variant calling, el ensamblaje y la investigación de enfermedades que vienen después. DeepConsensus alinea los subreads que observan varias veces la misma molécula de ADN en una representación bidimensional con gaps y, a partir de ella, predice una consensus sequence más precisa.

\subsection{De la secuenciación por consenso circular a la corrección aprendida}

Después de la página separadora de la sección, la página de la pregunta plantea directamente si Transformer puede mejorar la secuenciación de ADN. La página sobre el papel de DeepConsensus ubica al modelo dentro del pipeline de secuenciación: no genera secuencias biológicas nuevas, sino que aprovecha las observaciones repetidas que produce el instrumento para corregir los base call. La página de PacBio circular consensus sequencing (CCS) explica que la misma molécula circular puede leerse varias veces y que los múltiples subread aportan evidencia redundante.

\lecturefigure{slide-58-section-genomics.jpg}{Cambio de sección: aplicaciones genómicas}{Diapositiva V058 recuperada de la grabación oficial; Stanford CS25 Lecture 17}

\lecturefigure{slide-59-deepconsensus-question.jpg}{Pregunta: ¿puede Transformer mejorar la secuenciación de ADN?}{Diapositiva V059 recuperada de la grabación oficial; Stanford CS25 Lecture 17}

\lecturefigure{slide-60-deepconsensus-role.jpg}{Posición de DeepConsensus en el flujo de secuenciación}{Diapositiva V060 recuperada de la grabación oficial; Stanford CS25 Lecture 17}

\lecturefigure{slide-61-pacbio-consensus-sequencing.jpg}{Mecanismo de lectura repetida de PacBio circular consensus sequencing}{Diapositiva V061 recuperada de la grabación oficial; Stanford CS25 Lecture 17}

\begin{knowledgebox}{Concepto previo: read, subread y consensus}
Un \term{read} es un segmento de secuencia de bases que entrega el secuenciador; cuando la plantilla circular se recorre varias veces se producen múltiples \term{subread}; al alinearlos según la posición de referencia e integrarlos se obtiene el \term{consensus read}. Las observaciones repetidas permiten promediar los errores aleatorios, pero las insertion, las deletion y los sesgos sistemáticos exigen que el modelo entienda el alignment, en lugar de aplicar una votación por mayoría columna por columna.
\end{knowledgebox}

\subsection{Etiquetas de entrada y representación gap-aware}

La página de la tarea de corrección alinea varios subread con el draft actual, y la página de etiquetas indica que la salida es la consensus base verdadera. La dificultad está en que las inserciones y las deleciones producen gaps; si el modelo trata un gap como un espacio vacío común, las relaciones de posición se desordenan. Por eso DeepConsensus usa una gap-aware representation que codifica la identidad de la base, la calidad, la strand, la posición y los eventos de alineamiento como características que Transformer puede procesar.

\lecturefigure{slide-62-ccs-correction-task.jpg}{CCS correction: predecir la secuencia consenso a partir de varios noisy subread}{Diapositiva V062 recuperada de la grabación oficial; Stanford CS25 Lecture 17}

\lecturefigure{slide-63-deepconsensus-labels.jpg}{Etiquetas de entrenamiento: base correcta y posición de alineamiento}{Diapositiva V063 recuperada de la grabación oficial; Stanford CS25 Lecture 17}

\lecturefigure{slide-64-deepconsensus-transformer.jpg}{Arquitectura gap-aware Transformer de DeepConsensus}{Diapositiva V064 recuperada de la grabación oficial; Stanford CS25 Lecture 17}

\begin{lstlisting}[caption={Flujo de corrección al estilo DeepConsensus (pseudocódigo conceptual)}]
def correct_consensus(aligned_subreads):
    features = encode_bases_gaps_quality(aligned_subreads)
    hidden = gap_aware_transformer(features)
    base_logits, quality_logits = output_heads(hidden)
    sequence = decode_aligned_bases(base_logits)
    quality = calibrate_quality(quality_logits)
    return sequence, quality
\end{lstlisting}

\begin{warningbox}{Por qué un gap no es padding}
El padding significa «aquí no hay entrada», mientras que un alignment gap indica que, respecto de otra secuencia, ocurrió una insertion o una deletion: es un evento biológico y de medición que hay que explicar. Enmascarar un gap por error altera la correspondencia de posiciones posterior y hace que el modelo pierda información justamente en las regiones de indel, que son las más difíciles.
\end{warningbox}

Esta entrada también muestra que «las observaciones múltiples no son un lote común». Los distintos subread corresponden a la misma molécula física; el modelo debe integrar la evidencia en cada columna alineada y, al mismo tiempo, reconocer la baja calidad local de un pass determinado; si trata los subread como muestras independientes entre sí, pierde la información de consenso. En cambio, si se aplica directamente una votación por mayoría, no se pueden aprovechar la strand, las puntuaciones de calidad ni los errores dependientes del contexto. El papel de Transformer es precisamente reunir en cada posición alineada la evidencia de varios read y comprender el insertion/deletion pattern a través de las columnas vecinas. El aumento de datos también debe conservar esta estructura geométrica: no se pueden reordenar las columnas al azar ni combinar read de moléculas distintas.

\subsection{Alignment loss y base quality}

La token cross-entropy común supone que las posiciones predichas y las posiciones de las etiquetas se corresponden una a una, pero las insertion/deletion de las secuencias de secuenciación provocan desplazamientos globales. La página de alignment loss subraya que hace falta permitir un alineamiento local antes de penalizar las diferencias. En forma abstracta puede expresarse como
$$
\mathcal{L}_{\text{align}}(\hat Y,Y)=
\min_{a\in\mathcal{A}(\hat Y,Y)}
\sum_{(i,j)\in a} c(\hat y_i,y_j)+\lambda\,g(a),
$$
donde $\hat Y$ es la secuencia predicha, $Y$ es la truth, $\mathcal{A}$ es el conjunto de alignment permitidos, $c$ es el costo de un base mismatch, $g(a)$ es el costo de los gaps y $\lambda$ controla la penalización de los gaps. Lo que expresa la fórmula es que, si dos secuencias biológicas idénticas difieren solo en una inserción, no deben considerarse erróneamente equivocadas todas las posiciones siguientes.

\lecturefigure{slide-65-alignment-loss.jpg}{Alignment loss especial: manejo de inserciones, deleciones y desfases de posición}{Diapositiva V065 recuperada de la grabación oficial; Stanford CS25 Lecture 17}

El modelo también entrega la base quality, que expresa la probabilidad de error con una puntuación de estilo Phred:
$$
Q=-10\log_{10}p_{\mathrm{error}}.
$$
Donde $p_{\mathrm{error}}$ es la probabilidad de que esa base sea errónea; cuanto mayor es $Q$, más confiable es la predicción. Las puntuaciones de calidad deben estar calibradas: si se declara $Q=30$, la tasa de error entre predicciones similares debe acercarse a una en mil, y no basta con que el valor del logit sea grande.

\lecturefigure{slide-66-base-quality-output.jpg}{DeepConsensus entrega a la vez la base de consenso y la puntuación de calidad}{Diapositiva V066 recuperada de la grabación oficial; Stanford CS25 Lecture 17}

\begin{importantbox}{Dos cabezas de salida para dos problemas de decisión}
La base head responde «cuál es la secuencia» y la quality head responde «qué tan seguros estamos». El variant caller posterior necesita usar ambas; un modelo muy preciso pero con un exceso de confianza grave seguirá amplificando los errores. Por eso la calibración no es una métrica accesoria, sino el contrato de interfaz del pipeline de secuenciación.
\end{importantbox}

\subsection{Read-level accuracy e impacto en sistemas reales}

La página de read-level accuracy compara el modelo con los métodos de consensus existentes. La tendencia principal es que DeepConsensus aumenta la proporción de read de alta precisión, en lugar de obtener solo una ganancia mínima en el per-base error promedio. La página de real-world impact muestra además los beneficios en la etapa de despliegue: lecturas más precisas pueden reducir la secuenciación repetida necesaria para alcanzar la misma calidad, elevar el rendimiento y bajar el costo de los análisis posteriores.

\lecturefigure{slide-67-read-level-accuracy.jpg}{Mejora de DeepConsensus en read-level accuracy}{Diapositiva V067 recuperada de la grabación oficial; Stanford CS25 Lecture 17}

\lecturefigure{slide-68-real-world-impact.jpg}{Impacto en sistemas reales: mayor precisión, rendimiento y beneficios en costo}{Diapositiva V068 recuperada de la grabación oficial; Stanford CS25 Lecture 17}

\begin{knowledgebox}{Lectura de la figura: de las métricas del modelo a las métricas del sistema}
Primero hay que ver si baja el base error, luego la proporción de read completos que superan el umbral de alta calidad y, por último, cuántos datos útiles pueden entregarse con el mismo presupuesto. Las dos primeras capas corresponden al modelo y a la calidad de la secuenciación; solo la última refleja el beneficio para el sistema. Los beneficios del despliegue dependen también del hardware, la latencia del pipeline, las repeticiones por fallas y la compatibilidad con las herramientas posteriores, por lo que no pueden extrapolarse a partir de una sola exactitud.
\end{knowledgebox}

La evaluación del despliegue también debe considerar la distribución de los errores, y no solo la curva promedio. Algunas regiones repetitivas, homopolímeros o regiones con GC extremo pueden seguir siendo difíciles; si los errores se concentran en genes de relevancia clínica, la read accuracy promedio ocultará el riesgo posterior. El equipo de sistemas necesita hacer una paired comparison entre el pipeline anterior y el modelo nuevo con las mismas muestras y el mismo compute budget, y verificar si las tareas posteriores, como variant calling, assembly o methylation, mejoran de verdad. Solo cuando el beneficio del modelo se mantiene después de atravesar estas interfaces, el real-world impact deja de ser una conversión de métricas fuera de línea en cifras de mercado.

\teachervoice{El ponente presenta DeepConsensus como un ejemplo de algo que «ya tiene impacto en el mundo real»: Transformer no solo escribe texto o sube posiciones en un benchmark, sino que entra en la infraestructura de secuenciación y cambia la calidad y el costo de los datos que se pueden obtener. Este caso nos recuerda que el valor de la biomedical AI a menudo se amplifica a través de todo el pipeline.}

\subsection{Resumen de la sección}

DeepConsensus aprovecha múltiples observaciones noisy de la misma molécula de ADN y, mediante un gap-aware Transformer, una alignment-sensitive loss y la predicción de puntuaciones de calidad, genera consensus read más precisos. Su cadena completa de evidencia va de la base al read y de ahí al rendimiento y el costo del sistema de secuenciación, y muestra cómo las métricas del modelo se convierten en valor de infraestructura.
\section{Enformer: predicción de la regulación de largo alcance y de la expresión génica a partir de la secuencia de DNA}

DeepConsensus mejora primero la pregunta de «si la secuencia que se observa es correcta»; Enformer pregunta después «qué efecto regulador producirá esta secuencia». Muchas variantes asociadas con enfermedades se encuentran en la non-coding region: no alteran directamente la codificación de proteínas, pero pueden influir en la transcription factor binding, la enhancer--promoter interaction y la gene expression. Para predecir este efecto, el modelo necesita abarcar un contexto que va mucho más allá del motif local.

\subsection{Por qué los GWAS llevan el problema hacia las regiones no codificantes}

La diapositiva del artículo presenta Enformer, y la diapositiva de GWAS muestra que, a medida que crece la cohort, el número de disease-associated loci descubiertos aumenta con rapidez. La diapositiva Coding versus non-coding expone el hecho fundamental: una gran cantidad de señales de asociación se ubica en regiones que no codifican proteínas. Estas pueden influir en los genes mediante la regulación, en lugar de alterar directamente un amino acid; por ello, los flujos de trabajo tradicionales que solo examinan las coding mutation no pueden explicar la mayoría de las asociaciones.

\lecturefigure{slide-69-gene-expression-paper.jpg}{Trabajo asignado: Effective Gene Expression Prediction from Sequence}{Diapositiva V069 recuperada de la grabación oficial; Stanford CS25 Lecture 17}

\lecturefigure{slide-70-gwas-growth.jpg}{Crecimiento de la escala de los GWAS y descubrimiento de loci asociados}{Diapositiva V070 recuperada de la grabación oficial; Stanford CS25 Lecture 17}

\lecturefigure{slide-71-coding-vs-noncoding.jpg}{Diferencias funcionales entre variantes coding y non-coding}{Diapositiva V071 recuperada de la grabación oficial; Stanford CS25 Lecture 17}

\begin{knowledgebox}{Concepto previo: GWAS, variant y causalidad}
Un \term{GWAS} compara la asociación entre las genetic variant y el fenotipo en una población; una \term{variant} es una diferencia de bases de un individuo respecto del genoma de referencia; una \term{association} indica una correlación estadística, que no equivale automáticamente a una causal variant. El desequilibrio de ligamiento, la estructura poblacional y los factores ambientales pueden hacer que varios loci cercanos resulten significativos a la vez; por eso las predicciones del modelo se usan principalmente para establecer prioridades y formular hipótesis sobre mecanismos.
\end{knowledgebox}

\subsection{Enhancer, promoter y acción a distancia}

La diapositiva Transcription regulation muestra la relación entre promoter, enhancer, transcription factor y gene expression. La diapositiva Variant disrupts binding explica además que el cambio de una sola base no codificante puede alterar la unión de un factor y, en consecuencia, modificar la expresión de un gen lejano. La convolución local puede reconocer un motif, pero difícilmente conecta de forma directa un enhancer y un promoter separados por decenas de miles de bases; esa es precisamente la motivación para usar attention en el genoma.

\lecturefigure{slide-72-transcription-regulation.jpg}{Regulación de la transcripción: promoter, enhancer y transcription factor}{Diapositiva V072 recuperada de la grabación oficial; Stanford CS25 Lecture 17}

\lecturefigure{slide-73-variant-disrupts-binding.jpg}{Cómo una variante no codificante interrumpe la unión y modifica la expresión}{Diapositiva V073 recuperada de la grabación oficial; Stanford CS25 Lecture 17}

\begin{warningbox}{Predecir la regulación no es diagnosticar una enfermedad}
El modelo puede predecir que cierto allele probablemente modifique un track de chromatin o de expression, pero el fenotipo de una enfermedad también depende del tejido, la etapa del desarrollo, el ambiente, múltiples genes y el contexto clínico. Las puntuaciones de efecto de variantes sirven para seleccionar candidatos experimentales; no deben traducirse directamente en una probabilidad individual de padecer la enfermedad.
\end{warningbox}

\subsection{Entrada, salida y arquitectura híbrida de Enformer}

La diapositiva de la tarea toma una DNA sequence larga como entrada y predice numerosos genomic tracks; la diapositiva de Enformer muestra la combinación de convolution y Transformer. La convolución comprime primero los motif locales, la attention establece después interacciones entre posiciones en un rango más amplio y las cabezas de salida predicen señales para distintos tipos celulares y mediciones experimentales. El modelo lee un contexto de aproximadamente 200,000 bases, por lo que tiene más posibilidades que los modelos de ventana corta de conectar elementos reguladores distantes.

\lecturefigure{slide-74-dna-sequence-prediction-task.jpg}{Predicción de múltiples señales genómicas funcionales a partir de la secuencia de DNA}{Diapositiva V074 recuperada de la grabación oficial; Stanford CS25 Lecture 17}

\lecturefigure{slide-75-enformer-model.jpg}{Enformer: modelo de largo alcance que combina convolución y Transformer}{Diapositiva V075 recuperada de la grabación oficial; Stanford CS25 Lecture 17}

El modelo puede abstraerse como
$$
\hat Z=f_{\theta}(X),\qquad
X\in\{A,C,G,T\}^{L},\quad
\hat Z\in\mathbb{R}^{B\times K}.
$$
donde $X$ es la secuencia de DNA de longitud $L$, $B$ es el número de bins de posición de salida, $K$ es el número de assay, cell type o genomic track distintos, y $\hat Z$ es la señal predicha. El effect de una variante puede expresarse como la diferencia entre las dos secuencias de allele:
$$
\Delta_k=f_{\theta}(X_{\text{alt}})_k-f_{\theta}(X_{\text{ref}})_k.
$$
donde $X_{\text{ref}}$ y $X_{\text{alt}}$ difieren únicamente en la variant candidata, y $\Delta_k$ es el cambio predicho en el track $k$. Esta diferencia es una counterfactual estimate dentro del modelo, no una prueba experimental de causalidad.

La arquitectura híbrida también refleja una división del trabajo por escalas. La convolución tiene un sesgo inductivo fuerte hacia los motif locales y puede reconocer con eficiencia patrones cortos de unión de factores de transcripción; el pooling amplía gradualmente el campo receptivo y reduce la longitud de la secuencia; el Transformer permite después la interacción entre posiciones lejanas ya comprimidas. Aplicar attention a resolución completa sobre 200,000 bases desde el inicio sería muy costoso; si solo se apilaran convoluciones, la señal de un enhancer lejano tendría que atravesar muchas capas para llegar al promoter. Enformer no es «la attention que sustituye a la convolución», sino una asignación de los patrones locales y de la regulación global a los componentes adecuados.

\subsection{Thousands of tracks: la salida no es una sola etiqueta}

La diapositiva Tracks muestra que el modelo debe predecir simultáneamente una gran cantidad de señales experimentales, por ejemplo la chromatin accessibility, la transcription factor binding y la gene expression en distintos tipos celulares. La salida multitarea permite que la representación compartida absorba procesos biológicos relacionados, pero también aumenta el label noise y el desequilibrio de los datos: algunos track tienen datos abundantes, mientras que ciertos tejidos y poblaciones están claramente subrepresentados.

\lecturefigure{slide-76-enformer-tracks.jpg}{Genomic tracks de múltiples tipos celulares predichos por Enformer}{Diapositiva V076 recuperada de la grabación oficial; Stanford CS25 Lecture 17}

\begin{knowledgebox}{Lectura de la figura: entender un genomic track como «función de la posición»}
El eje horizontal es la posición en el genoma y el eje vertical es la intensidad de cierta señal experimental; los distintos colores o filas representan assay, cell type o valores del modelo frente a valores reales. Primero hay que ver si las posiciones de los picos coinciden, luego si las intensidades concuerdan y, por último, verificar si un cambio en un enhancer lejano se corresponde con una modificación de la expression cerca del promoter. Mirar solo el coeficiente de correlación global puede ocultar picos raros pero decisivos.
\end{knowledgebox}

\subsection{Resultados y límites de uso}

La diapositiva de resultados compara las predicciones relacionadas con promoter y enhancer, y muestra que un contexto más largo mejora la regulatory activity prediction. La diapositiva Application takeaways resume los usos: explicar la non-coding variation, predecir la gene expression y apoyar la variant prioritization. El valor científico del modelo consiste en reducir el espacio de búsqueda experimental: identificar, entre decenas de miles de loci asociados, los candidatos que más merecen un assay.

\lecturefigure{slide-77-promoter-enhancer-results.jpg}{Mejora de Enformer en la predicción de la promoter y enhancer activity}{Diapositiva V077 recuperada de la grabación oficial; Stanford CS25 Lecture 17}

\lecturefigure{slide-78-application-takeaways.jpg}{Conclusiones de aplicación y limitaciones de Enformer}{Diapositiva V078 recuperada de la grabación oficial; Stanford CS25 Lecture 17}

\teachervoice{El ponente establece primero la necesidad con «la gran mayoría de los GWAS hits están en la non-coding region» y después explica la enhancer--promoter interaction a larga distancia. Este orden es importante: la attention no se usa porque esté de moda, sino porque el mecanismo biológico plantea un problema que excede el receptive field local.}

\begin{importantbox}{La cadena de evidencia correcta de Enformer}
Un contexto largo de DNA mejora la predicción de múltiples genomic track y de la expression; la diferencia entre el allele reference/alternate permite ordenar las variant; los experimentos verifican después la unión, la expresión y el fenotipo. El modelo vuelve más inteligente la decisión de «dónde experimentar», pero no elimina los experimentos ni convierte automáticamente una statistical association en un disease mechanism.
\end{importantbox}

Si se quiere usar Enformer para la variant prioritization, también hay que elegir los track que correspondan al tejido de la enfermedad. En una enfermedad cardiaca, examinar solo la puntuación global de una línea celular cualquiera puede producir candidatos biológicamente irrelevantes; en variantes poblacionales, además, hay que comprobar si los datos de entrenamiento representan la ancestry objetivo. Un informe más sólido enumera la puntuación del modelo, los cell type relevantes, los gene cercanos, la evidencia eQTL/epigenomic existente y la viabilidad experimental, para que los investigadores conozcan el fundamento del orden. Aquí el modelo desempeña el papel de un evidence aggregator, no de un juez final.

\subsection{Resumen de la sección}

Enformer transforma secuencias largas de DNA en genomic tracks de múltiples tipos celulares y múltiples assay; usa convoluciones para extraer motif locales y un Transformer para establecer relaciones reguladoras de largo alcance. Puede mejorar la predicción de la gene-expression y del variant-effect, y ofrece una base para ordenar los experimentos sobre variantes no codificantes; sin embargo, asociación, predicción y causalidad siguen siendo tres niveles distintos de evidencia.
\section{Síntesis y ampliación}

Al final, la clase pasa de los cuatro casos a la foundation biomedical AI. El juicio central del ponente no es «si» habrá un cambio, sino «cuándo y de qué manera» ocurrirá. Para entender esta frase hay que condensar lo anterior en una cadena completa: los modelos clínicos procesan el conocimiento y la comunicación humanos, los modelos de proteínas conectan sequence con function language, DeepConsensus mejora la calidad de la observación y Enformer explica el contexto regulatorio. La oportunidad de los sistemas futuros está en que estas capas se aporten evidencia entre sí, no en meter todos los datos en una interfaz de chat.

\subsection{Un marco común extraído de los cuatro modelos}

La página de apertura de la sección sobre el futuro marca primero el cambio de nivel de la discusión y luego plantea el «when and how». No es una predicción temporal, sino un problema de ingeniería y de ciencia: ¿qué datos pueden combinarse de forma legal y segura? ¿Qué salidas pueden verificarse con una metric automática y cuáles deben pasar a un clinician o a un wet lab? ¿Cómo se diseñan los beneficios de escala de un foundation model multimodal junto con la provenance, la calibration y la selective prediction?

\lecturefigure{slide-79-section-future.jpg}{Cambio de sección: el futuro de la Biomedical AI}{Diapositiva V079 recuperada del video oficial; Stanford CS25 Lecture 17}

\lecturefigure{slide-80-future-how-not-if.jpg}{La pregunta sobre el futuro: no si ocurrirá, sino cuándo y cómo}{Diapositiva V080 recuperada del video oficial; Stanford CS25 Lecture 17}

\begin{longtable}{@{}L{0.19\textwidth}L{0.25\textwidth}L{0.25\textwidth}L{0.19\textwidth}@{}}
\toprule
Caso & Entrada y reto & Mecanismo del modelo & Validación externa \\
\midrule
Med-PaLM & Preguntas médicas; conocimiento, comunicación y riesgo & Scaling, instruction tuning, prompt tuning, abstention & Evaluación multidimensional por médicos y usuarios no especialistas \\
Performer / Protein LM & Secuencias de aminoácidos muy largas & kernel attention y aprendizaje de representaciones & Estructura, bases de datos y experimentos \\
ProtNLM & De secuencia de proteína a texto libre & encoder--decoder, T5 multitarea & Provenance de las anotaciones, expertos y experimentos \\
DeepConsensus & noisy subreads alineados & gap-aware attention, alignment loss, quality head & base/read accuracy y rendimiento del sistema \\
Enformer & De DNA largo a genomic tracks & convolution + long-range attention & Trazas experimentales, expresión y variant assay \\
\bottomrule
\end{longtable}

\begin{importantbox}{Principio unificador: la salida del modelo debe reconectarse con su mundo correspondiente}
Las respuestas de texto se reconectan con el consenso clínico y con los usuarios; las descripciones de proteínas, con las bases de datos y los experimentos; la corrección de errores de secuenciación, con los read y el pipeline; las predicciones regulatorias, con los assay y el fenotipo. Un foundation model puede compartir representaciones y datos, pero no puede compartir un criterio de evaluación difuso de «parece razonable».
\end{importantbox}

Esta tabla resumen muestra además que las partes compartidas y las partes especializadas de un foundation model deben diseñarse por separado. Un encoder compartido puede aprender representaciones intermodales, recuperación unificada y generación condicionada; los output head, la loss, la calibration y los reviewer especializados se encargan de restringir la salida a la semántica clínica, proteica o genómica. Si en nombre de la «unificación» se elimina la validación especializada, el sistema será más sencillo en la interfaz, pero más frágil en la evidencia. Por el contrario, la modularidad no impide el entrenamiento de extremo a extremo; solo exige que cada interfaz crítica pueda auditarse por separado: de dónde proviene la entrada, cómo se transmite la representación, quién confirma la salida y cómo se escala una falla.

\subsection{Foundation biomedical AI: integrar sin borrar las diferencias}

La página de foundation biomedical AI expresa la visión de integración con un disco multimodal. La página siguiente coloca lado a lado a un physician y a un scientist célebres; con ello, el ponente explica que la IA del futuro no tiene por qué separar la clinical application de la biological application: la observación clínica puede plantear preguntas sobre mecanismos, los modelos moleculares pueden generar candidatos y los resultados experimentales actualizan a su vez el conocimiento clínico. Aquí, «foundation» se parece más a una interfaz entre capas que a una simple cantidad mayor de parámetros.

\lecturefigure{slide-81-foundation-biomedical-ai.jpg}{La visión de integración multimodal de la foundation biomedical AI}{Diapositiva V081 recuperada del video oficial; Stanford CS25 Lecture 17}

\lecturefigure{slide-82-biomedical-ai-scientists.jpg}{Los papeles del médico y del científico pueden conectarse dentro de un sistema de IA}{Diapositiva V082 recuperada del video oficial; Stanford CS25 Lecture 17}

\teachervoice{El ponente dice que las clinical applications y las biological applications no tienen por qué permanecer separadas para siempre; solo al conectarlas pueden surgir nuevos insight, acelerarse la biomedical research y, finalmente, volver a los problemas de enfermedad y salud. Este juicio debe escribirse como una línea de investigación, no como la afirmación de que en 2023 ya existía un ciclo cerrado unificado.}

\begin{warningbox}{Las tres confusiones más frecuentes en la integración entre capas}
Primera: tomar una correlación clínica por causalidad molecular. Segunda: tomar la confidence del modelo por reproducibilidad experimental. Tercera: tomar por generalización la filtración del mismo paciente o de la misma base de datos tanto al entrenamiento como a la prueba. Cuanto más se avanza hacia lo foundation, más importantes son la provenance de los datos, la partición temporal, la privacidad y la validación independiente.
\end{warningbox}

\subsection{¿En qué campo obtendrá la IA su primer Nobel Prize?}

La penúltima página deja la pregunta a los estudiantes: ¿en qué campo propiciará primero la IA un Nobel-level discovery? No se trata de adivinar el premio, sino de preguntar «qué evidencia cuenta como avance científico». Si la IA solo condensa la literatura existente, su contribución es la de una herramienta; solo se acerca a un descubrimiento si propone mecanismos nuevos, orienta experimentos y es reproducida de manera independiente. Lo particular de la biomedicina es que tiene a la vez un enorme espacio sin resolver, experimentos medibles y un valor humano directo, y también los umbrales éticos y de validación más altos.

\lecturefigure{slide-83-ai-nobel-question.jpg}{Pregunta abierta: ¿en qué campo obtendrá la IA su primer Nobel Prize?}{Diapositiva V083 recuperada del video oficial; Stanford CS25 Lecture 17}

\lecturefigure{slide-84-thank-you.jpg}{Cierre de la clase y agradecimiento a los colaboradores}{Diapositiva V084 recuperada del video oficial; Stanford CS25 Lecture 17}

\teachervoice{El ponente deja la pregunta en el aula sin dar un año ni una respuesta definitiva; después comenta que espera que la mayoría piense en medicine. Lo que de verdad vale la pena conservar es la escala de evaluación: para que la IA tenga un impacto Nobel-level en medicina, debe pasar de la predicción al descubrimiento reproducible, y del descubrimiento al beneficio en salud.}

\subsection{Siete puntos de verificación para la práctica de ingeniería}

\begin{enumerate}
  \item \textbf{Definir primero la unidad de secuencia.} La semántica posicional de token, event, amino acid y nucleotide no debe mezclarse.
  \item \textbf{Escribir primero la validación externa y después elegir el modelo.} Si la salida no puede reconectarse con un clinician, una database, un assay o un pipeline, difícilmente un benchmark representará su valor real.
  \item \textbf{Hacer explícito el costo del contexto largo.} Registrar la sequence length, la complejidad de la attention, la estrategia de truncamiento y las relaciones de largo alcance que podrían perderse.
  \item \textbf{Convertir la incertidumbre en una interfaz.} La selective prediction, el quality score y la calibration deben servir a una ruta posterior de escalamiento, no limitarse a trazar una curva.
  \item \textbf{Conservar los riesgos multidimensionales.} La accuracy, el missing content, el harm, el bias y la helpfulness no deben reducirse prematuramente a una sola puntuación.
  \item \textbf{Distinguir predicción, asociación y causalidad.} Un attention map, un variant score y una anotación generada pueden plantear hipótesis, pero no sustituyen a los experimentos.
  \item \textbf{Respetar los límites históricos.} Responder las preguntas de la clase con los modelos y la evidencia de ese momento; los avances posteriores deben ir en una actualización aparte, no reescribir la cadena causal de los apuntes originales.
\end{enumerate}

\begin{importantbox}{Condensación final}
Lo esencial del Biomedical Transformer no es «convertir en lenguaje» la medicina y la biología, sino usar una herramienta unificada de modelado de contexto para conectar secuencias distintas, conservando para cada dominio su representación, pérdida, calibración y validación propias. Un foundation system genuino no es un modelo que responde a todo, sino un conjunto de interfaces confiables capaces de transmitir información entre varias capas de evidencia, de saber cuándo abstenerse de responder y de devolver los juicios críticos a los médicos y a los experimentos.
\end{importantbox}

\subsection{Lecturas adicionales}

\begin{itemize}
  \item Singhal et al., \href{https://www.nature.com/articles/s41586-023-06291-2}{\emph{Large Language Models Encode Clinical Knowledge} (Nature)}
  \item Choromanski et al., \href{https://arxiv.org/abs/2009.14794}{\emph{Rethinking Attention with Performers} (arXiv)}
  \item Google Research, \href{https://github.com/google-research/google-research/tree/master/protnlm}{\emph{ProtNLM}: código y documentación oficiales}
  \item Baid et al., \href{https://research.google/pubs/deepconsensus-improves-the-accuracy-of-sequences-with-a-gap-aware-sequence-transformer/}{\emph{DeepConsensus} (Google Research)}
  \item Avsec et al., \href{https://www.nature.com/articles/s41592-021-01252-x}{\emph{Effective Gene Expression Prediction from Sequence by Integrating Long-range Interactions} (Nature Methods)}
\end{itemize}

Estos cinco trabajos representan, respectivamente, la evaluación clínica, la attention para secuencias largas, la interfaz de lenguaje científico, la infraestructura de secuenciación y la genómica regulatoria. Se recomienda leerlos aplicando de manera uniforme las cinco preguntas de esta clase: cómo se representa la entrada, por qué el contexto es difícil, de dónde viene la supervisión, cómo se valida la salida y qué no se puede concluir. Así, los resultados de los artículos se convierten en métodos de investigación transferibles y no en una lista de nombres de modelos.

\end{document}
